% New records for equal sums of like powers via lattice sieving
% Build: latexmk -pdf main.tex
% All tables and the numbers set by macros (\NumImproved etc.) are generated from the
% verified result files by tools/make_tables.py, tools/controls.py, tools/ablation_table.py,
% tools/nsweep_table.py and tools/logstats.py.
\documentclass[11pt]{amsart}
\usepackage[utf8]{inputenc}
\usepackage{amssymb,amsmath}
\usepackage{cite}
\usepackage{booktabs,longtable,array}
\usepackage{algorithm,algpseudocode}
\usepackage{pdflscape,capt-of}
\usepackage{xcolor}
\usepackage{microtype}
\usepackage[hidelinks]{hyperref}

% generated by tools/make_tables.py -- do not edit
\newcommand{\NumImproved}{197}
\newcommand{\NumTies}{9}
\newcommand{\TermsSaved}{3,266}
\newcommand{\MaxGain}{58}
\newcommand{\MaxGainCell}{$(30,3)$}
\newcommand{\MaxGainOld}{189}
\newcommand{\MaxGainNew}{131}
\newcommand{\MaxBase}{115}
\newcommand{\NumChecked}{86,470}
\newcommand{\NumFullK}{eight}
\newcommand{\FullKList}{22, 23, 24, 25, 27, 28, 29, 30}
\newcommand{\KMinImproved}{17}
\newcommand{\KMaxImproved}{30}
\newcommand{\NumCells}{221}
\newcommand{\NumNotImproved}{24}
\newcommand{\NumDerived}{41}
\newcommand{\NumBeatIdent}{194}
\newcommand{\NoIdentCells}{$(20,8)$, $(21,7)$ and $(22,13)$}
\newcommand{\NumTiesDerived}{1}
\newcommand{\TiesDerivedSentence}{One of the matched entries is also a derived bound, and for it our identity is the first of its type.}
\newcommand{\ClosureExtra}{$(21,14)$}
\newcommand{\NumClosureExtra}{one}
\newcommand{\NumClosureLower}{10}
\newcommand{\NumCellsDerived}{43}
\newcommand{\NumCellsIdent}{178}
\newcommand{\NumImprovedIdent}{156}
\newcommand{\TermsSavedIdent}{2,402}
\newcommand{\TermsSavedDerived}{864}
\newcommand{\HNumCells}{1,088}
\newcommand{\HNumCellsDerived}{475}
\newcommand{\HNumCellsIdent}{613}
\newcommand{\HNumImproved}{1,071}
\newcommand{\HNumDerived}{465}
\newcommand{\HNumImprovedIdent}{606}
\newcommand{\HTermsSaved}{2,218,809}
\newcommand{\HMaxGain}{25,340}
\newcommand{\HMaxGainCell}{$(99,1)$}
\newcommand{\HMaxGainOld}{30,611}
\newcommand{\HMaxGainNew}{5,271}
\newcommand{\HNumK}{68}
\newcommand{\HNumFullK}{65}
\newcommand{\HNoneK}{}
\newcommand{\HNoneSentence}{Every power has at least one improved entry.}
\newcommand{\HMaxBase}{112,640}
\newcommand{\AllImproved}{1,300}
\newcommand{\KThirtyOneImproved}{16}
\newcommand{\KThirtyOneSaved}{864}
\newcommand{\KThirtyTwoImproved}{16}
\newcommand{\KThirtyTwoSaved}{255}

% generated by tools/ablation_table.py -- do not edit
\newcommand{\AblK}{20, 22, 24, 26, 28, 30}
\newcommand{\AblSeeds}{three}
\newcommand{\StepsMax}{16}
\newcommand{\StepsMedian}{4}
\newcommand{\StepsAtCap}{0}
\newcommand{\StepsTasks}{152,460}
\newcommand{\SatSieves}{426}
\newcommand{\SatFailed}{0}
\newcommand{\SatFinalFailed}{0}
\newcommand{\SatRuns}{36}
\newcommand{\AblImpA}{1}
\newcommand{\AblImpB}{0}
\newcommand{\AblImpC}{14}
\newcommand{\AblImpD}{17}
\newcommand{\AblImpE}{12}

% generated by tools/logstats.py -- do not edit
\newcommand{\CoreHoursNRP}{634}
\newcommand{\CoreHoursTuring}{2,020}
\newcommand{\CoreHoursTotal}{2,654}
\newcommand{\GPUHours}{678}

% generated by tools/congruence.py -- do not edit
\newcommand{\ForcedK}{26}
\newcommand{\ForcedSieveImp}{91}
\newcommand{\ForcedAllImp}{399}
\newcommand{\ForcedCells}{416}
\newcommand{\UnforcedK}{42}
\newcommand{\UnforcedSieveImp}{671}
\newcommand{\UnforcedCells}{672}
\newcommand{\ForcedMainK}{17, 19, 25, 31, 32}

% generated by tools/controls.py -- do not edit
\newcommand{\CtrlRImp}{191}
\newcommand{\CtrlRSaved}{3,173}


\newtheorem{proposition}{Proposition}[section]
\newtheorem{lemma}[proposition]{Lemma}
\theoremstyle{definition}
\newtheorem{definition}[proposition]{Definition}
\theoremstyle{remark}
\newtheorem{remark}[proposition]{Remark}

\newcommand{\Z}{\mathbb{Z}}
\newcommand{\GH}{\mathrm{GH}}

\begin{document}

\title[New records for equal sums of like powers]{New records for equal sums of like powers\\ via lattice sieving}

\author{J. G. Childers}
\address{Department of Physics, California State University, Fullerton, 800 N. State College Blvd., Fullerton, CA 92831, USA}
\email{gchilders@fullerton.edu}

\subjclass[2020]{Primary 11D41, 11Y50; Secondary 11H06, 11Y16}
\keywords{Equal sums of like powers, Euler's conjecture, integer relations, lattice reduction, lattice sieving}

\begin{abstract}
An equal sum of like powers of type $(k,m,n)$ is an identity
$a_1^k+\dots+a_m^k=b_1^k+\dots+b_n^k$ in positive integers; for each $k$ and $m$ one seeks
the least $n$. Between 2000 and 2003 many records for large $k$ were found by lattice reduction
of the integer relations among $1^k,\dots,N^k$. We return to this lattice with the General
Sieve Kernel and use the whole sieve database as the move set of a greedy descent on the
numbers of terms. With bases at most $120$ this improves \NumImproved{} of the \NumCells{}
entries with $m\le16$ and $17\le k\le30$ in Meyrignac's table of records; together with
completions of congruence solutions, which use larger bases, it improves \HNumImproved{} of the
\HNumCells{} entries for $33\le k\le100$; for 25 powers it also lowers the least known total number of terms. The
powers that resist the search are mostly those for which $k-j$, for a small $j\ge1$, has many
divisors $\ell-1$ with $\ell$ prime. Then the moment $\sum_x c_xx^j$, where $c_x$ is the signed
multiplicity of the base $x$, vanishes on almost all short
relations, a descent through such relations cannot change it, and the sublattice where it
vanishes appears to contain few relations with one short side; for most of these powers,
starting in other cosets removes this obstruction. All identities are verified in exact arithmetic,
and a standalone verifier is provided with the data.
\end{abstract}

\maketitle

%======================================================================
\section{Introduction}
%======================================================================

Euler conjectured that for $k\ge3$ a $k$-th power is never a sum of fewer than $k$ positive
$k$-th powers. Lander and Parkin disproved this for $k=5$ with
$27^5+84^5+110^5+133^5=144^5$ \cite{LanderParkin1966,LanderParkin1967}. Elkies
\cite{Elkies1988} disproved it for $k=4$, and Frye \cite{Frye1988} found the smallest
counterexample. Lander, Parkin and
Selfridge \cite{LanderParkinSelfridge1967} surveyed equal sums of like powers more generally
and conjectured that a nontrivial identity of type $(k,m,n)$ requires $m+n\ge k$. Ekl
\cite{Ekl1996,Ekl1998} found further identities, and Gerbicz, Meyrignac and Beckert
\cite{GerbiczMeyrignacBeckert2011} attribute to Ekl \cite{Ekl1998} the name ``Euler's
extended conjecture'' for this statement.

For large $k$ nothing beyond computation is known. For each $k$ and $m$ one seeks an identity
of type $(k,m,n)$ with $n$ as small as possible, and since 1999 the best known values have been
collected by J.-C.~Meyrignac \cite{MeyrignacEuler}. About three quarters of the dated entries
with $13\le k\le32$ in his table come from 2000--2003, when the problem was treated as an
integer relation problem and attacked with lattice basis reduction
\cite{LLL1982,SchnorrEuchner1994,ChildersSumsPage2000,CrumpEqualSums2001}. In the table no
entry with $k\ge11$ is dated later than 2013.

Lattice reduction has advanced considerably since then. BKZ~2.0
\cite{ChenNguyen2011} and the floating-point LLL algorithm of Nguyen and Stehl\'e
\cite{NguyenStehle2009} made strong reduction routine in fplll \cite{fplll}. Lattice sieving
\cite{NguyenVidick2008,MicciancioVoulgaris2010,BeckerDucasGamaLaarhoven2016} has overtaken
enumeration in practice, and the General Sieve Kernel (G6K) \cite{AlbrechtEtAl2019}
organized sieving around a persistent database of short vectors. Its GPU version solved
Darmstadt SVP challenges up to dimension $180$ on a single server with four GPUs
\cite{DucasStevensVanWoerden2021}. Lattice reduction has also been used in other Diophantine
searches, for rational points near curves and small values of $|x^3-y^2|$
\cite{Elkies2000} and for $abc$ triples \cite{Dokchitser2004}, and integer relation
algorithms such as PSLQ \cite{FergusonBaileyArno1999,BaileyBroadhurst2001} are a standard tool
of experimental mathematics. We know of no refereed account of lattice methods for equal sums
of like powers.

The paper has two aims. The first is to record the method of 2000--2003 while the programs
and records survive (Section~\ref{sec:history}). The second is to show that the same lattice,
searched with current tools, yields many new records, and to determine which parts of the
search are responsible. The components are not new individually: the relation lattice
dates from 2000, the weighted greedy objective from the program \texttt{ResultLC} of 2001, and the
sieve is G6K. The new element is the way they are combined. The descent
(Section~\ref{sec:descent}) uses the whole sieve database, up to $6\times10^7$
relations, as its set of moves, where the programs of 2000 used the $N-1$ rows of one
reduced basis, and it records every relation it examines, not only the local minima it
reaches.

The results (Section~\ref{sec:results}) improve \NumImproved{} of the \NumCells{} entries with
$m\le16$ and $17\le k\le30$, including every entry for $k\in\{\FullKList\}$.
Section~\ref{sec:experiments} separates the contributions of the stages. Our combination
procedure, applied to the identities already known, improves no entry of the table, while the
identities found by individual runs, without combination and without the known identities,
improve \CtrlRImp{} of the \NumImproved{} entries. The database alone reaches almost none of the
records, and the descent does. With the same descent, an equal number of rows from thousands of
independent BKZ reductions with block size 40 serves about as well as sieve vectors, and neither
move set is uniformly better; the advantage of the sieve is that it produces many more moves in
the same time.

Section~\ref{sec:high} turns to the companion table for $33\le k\le100$. There a congruence
obstruction (Proposition~\ref{prop:moment}) accounts for most of the powers that resist the
search; starting the descent in other cosets of a sublattice of the relation lattice, and
completing congruence solutions as in 2002, overcomes it for most of them, and \HNumImproved{}
of the \HNumCells{} entries are improved. In the two tables together we lower the least known
total $m+n$ for 25 powers (10 in the main table and 15 in the second). The identities, the
reference version of the tables and a verification program that needs only a Python interpreter
are in the data archive.

%======================================================================
\section{The problem and the record table}\label{sec:problem}
%======================================================================

\begin{definition}
Let $k\ge2$. An identity
\[
a_1^k+\dots+a_m^k=b_1^k+\dots+b_n^k,\qquad a_i,b_j\in\Z_{>0},
\]
in which no base occurs on both sides is of type $(k,m,n)$; we take $m\le n$. Bases may
repeat, and we write a repeated base with a multiplicity, as in $3\cdot7^k$. The identity is
\emph{primitive} if the bases have no common divisor greater than $1$.
\end{definition}

Equivalently, an identity is a nonzero vector $c\in\Z^N$, where $N$ is at least the largest
base, with
\begin{equation}\label{eq:relation}
\sum_{x=1}^{N} c_x\, x^k = 0.
\end{equation}
The positive and negative entries of $c$ give the two sides, and $|c_x|$ is the multiplicity
of $x$. Write $p(c)=\sum_{c_x>0}c_x$ and $q(c)=\sum_{c_x<0}|c_x|$. Then $c$ gives an identity
of type $(k,\min(p,q),\max(p,q))$ with $T=p+q=\|c\|_1$ terms, where
$\|c\|_1=\sum_x|c_x|$. We call $c$ nearly balanced if $\min(p,q)$ is close to $T/2$, and
unbalanced if $\min(p,q)$ is much smaller; the records with small $m$ are unbalanced. We write
$\|c\|_2=(\sum_xc_x^2)^{1/2}$ for the Euclidean length of $c$, which is what lattice reduction
and sieving make small; a norm without subscript, as in $\|c\|^2$, is the Euclidean norm.

Two operations produce an identity from another without new information. Multiplying every
base by an integer greater than $1$ preserves the type, and multiplying $c$ by an integer $t>1$ gives the type
$(k,tm,tn)$. We count only primitive identities, and we check that none of our identities
has a common factor in its multiplicities. Adding the same base to both sides is not an
operation in this sense, since the result violates the definition.

Meyrignac's table \cite{MeyrignacEuler} lists, for $2\le k\le32$ and $1\le m\le16$, the least
known $n$ for each $(k,m)$; a second table covers $33\le k\le100$. The entries are of two
kinds. Most are attained by an identity in the accompanying database, and a few of these are
proved or conjectured minimal. The others, which we call derived bounds, are obtained from a
neighboring entry by the rule that an identity of type $(k,m,n)$ gives the bound $n+i$ for the
entry $(k,m+i)$. Such a bound does not correspond to an identity in the sense of the
definition. Of the \NumCells{} entries with $17\le k\le30$ and $m\le16$, \NumCellsIdent{} are of
the first kind and \NumCellsDerived{} of the second; the first kind includes $(18,15)$, which is
empty in the table but attained by an identity of type $(18,15,15)$ in the database. Our counts
refer to this range, which our main searches cover; the entries for $k=31$ and $32$, which were
added later and needed partly different methods, are reported separately
(Section~\ref{sec:notimproved}), and a search for $k=16$ improved no entry. We compare our
results with the table and its database as accessed on 27 October 2026; the database was last
modified in November 2022. The reference value of an entry is the smaller of the table entry
and the least $n$ of an identity of type $(k,m,n)$ in the database. The archive records SHA-256
checksums of the files from which the reference values were extracted, and lists the reference
value of every entry.

%======================================================================
\section{The integer-relation searches, 2000--2003}\label{sec:history}
%======================================================================

The record table grew out of distributed enumeration. Meyrignac's search began in January 1999
and became the EulerNet client--server project in July 2000, with the server hosted by the
author for most of its life; from 2010 to 2011 it continued as a subproject of the BOINC project
yoyo@home (Appendix~\ref{app:chronology}). Exhaustive enumeration of this kind settles small $k$ up to a bound on the terms; for large
$k$ it is impractical, and lattice methods took its place.

\subsection{Integer relations}
Joe McLean formulated the search as $z^k=\sum_x c_x x^k$. In September 2000 the author
treated \eqref{eq:relation} as an integer relation problem among $1^k,\dots,N^k$, tested the
approach with the LLL algorithm \cite{LLL1982} of Mathematica before sharing it, and held the
records for $k\ge23$ on 24 and 25 September \cite{MeyrignacEuler}. Joe K.~Crump then
used the BKZ reduction of NTL \cite{SchnorrEuchner1994,NTL}, and from 26 September 2000 held
every record above $k=12$ \cite{CrumpEqualSums2001}. The author's program, first written in PARI/GP and
then as \texttt{SumsBKZ} using NTL's BKZ, reached version~1.3.2 by December 2000
\cite{ChildersSumsPage2000}. Meyrignac's site credits the author with the application of integer relations
to the problem \cite{MeyrignacEuler}.

\texttt{SumsBKZ} reduces the $N\times(N+1)$ basis
\[
B=\begin{pmatrix} I_N \;\big|\; \lambda\,(x_1^k,\dots,x_N^k)^{\mathsf T}\end{pmatrix},\qquad \lambda=10^4,
\]
with NTL's \texttt{BKZ\_FP} ($\delta=0.9999$), where the bases $x_i$ are $1,\dots,N$ in random
order or, for the table of prime bases, the first $N$ primes. Since $\lambda$ is large, the
short vectors of this lattice have last coordinate zero, and the reduced rows with last
coordinate zero are relations \eqref{eq:relation}. Each pass reduces one random ordering. Options restrict the reduction to
a random subset of the bases, always include the smallest bases, or, for even $k$, use $1$
and multiples of chosen primes. The last option addresses the following restriction. If $\ell$
is prime and $(\ell-1)\mid k$, then $x^k\equiv1\pmod\ell$ whenever $\ell\nmid x$, so every
relation satisfies $\sum_{\ell\nmid x}c_x\equiv0\pmod\ell$. The numbers of terms prime to $\ell$
on the two sides are then congruent modulo $\ell$, and using multiples of $\ell$ as bases removes
this restriction.

\subsection{Post-processing and combination}
BKZ produces vectors that are short in the Euclidean norm, whereas a record is judged by
$\min(p,q)$ and then by $p+q$. Fumitaka Yura's post-processing step, adopted in
\texttt{SumsBKZ} in December 2000, partly addresses this. Starting from each reduced
relation, it adds or subtracts another reduced relation whenever this lowers $\min(p,q)$, and
then repeats the process for $p+q$. Two further programs combined identities found in
different runs. \texttt{ResultLC}, whose documentation states that it uses the algorithm of Yura's
program \texttt{ysum}, applies the same kind of step to the objective
$w\min(p,q)+p+q$ for an increasing sequence of weights $w$. \texttt{NewLC}, written by Meyrignac in
January 2001, is a multi-pass beam search over sums and differences of known identities that
keeps the best bound for each $m$. Later variants of 2002 included \texttt{SumBKZ2} and a version
of \texttt{SumsBKZ} for large exponents.

\subsection{Congruence methods and completion}
Jaros\l aw Wr\'oblewski's searches, which produced many records for $9\le k\le20$, restrict
the bases to residue classes modulo prime powers before searching
\cite{WroblewskiProblemSet}; the method is otherwise unpublished. In 2002 he also supplied
congruence solutions for larger powers, and a variant of \texttt{SumsBKZ} completed them; the
records for $k=31$ and for several powers between $37$ and $97$ were found in this way. Given
multisets of bases $P$ and $Q$ with $\sum_{y\in P}y^k-\sum_{y\in Q}y^k=M^k\Delta$, this variant
appends the row $(0,\dots,0,\lambda\Delta)$ to $B$ and keeps the reduced vectors with last
coordinate zero that use this row with coefficient $\pm1$; up to sign they satisfy
$\sum_xc_xx^k=\Delta$, and the bases $Mx$ with multiplicities $c_x$ complete $P$ and $Q$ to an
identity. This is a form of Kannan's embedding of a closest-vector problem into a
shortest-vector problem \cite{Kannan1987}.

\subsection{The archive and a baseline}
The surviving archive of 2001--2003, which we call the 2003 archive, contains 1,042 identities
for $k\le32$, 3,384 for $33\le k\le100$ and 866 for the table of prime bases, all of which
verify. We rebuilt \texttt{SumsBKZ}~1.3.3, the last version that survives, against NTL~11.6 and
ported it to Python with fpylll \cite{fpylll}. With identical settings (10 passes,
$N=40,\dots,50$, $9\le k\le32$, block size~20) the original reproduced two of the archived
records that are reachable with bases at most $50$, and the port four, using 146\,s and 197\,s
of CPU time, respectively. This sample is small; it shows continuity of behavior, not
equivalence of the two implementations. The port is baseline~A of
Section~\ref{sec:experiments}.

%======================================================================
\section{Method}\label{sec:method}
%======================================================================

\subsection{The relation lattice}
For $k$ and $N$ let $v=(1^k,2^k,\dots,N^k)$ and
\[
L=L_{k,N}=\{\,c\in\Z^N : \langle c,v\rangle=0\,\}.
\]
This is a lattice of rank $d=N-1$, and since $\gcd(v)=1$ its determinant is $\|v\|$. Since
$1^k=1$, the vectors $e_x-x^ke_1$ ($2\le x\le N$) form a basis of $L$, but a poorly reduced
one. We obtain a reduced basis from one LLL reduction of the embedding of
Section~\ref{sec:history}: when exactly $N-1$ of the reduced rows have last coordinate zero, as
we check, these rows form a basis of $L$, since a lattice vector with last coordinate zero has
coefficient zero on the remaining row. We then reduce it further with BKZ of block size 20
(BKZ-20; similarly BKZ-30 and BKZ-40). All later computation takes place
in $L$, whose reduced basis has small entries.

\subsection{The Gaussian heuristic and the number of terms}\label{sec:gh}
Let
\begin{equation}\label{eq:gh}
\GH(L)=\Bigl(\frac{d}{2\pi e}\Bigr)^{1/2}\|v\|^{1/d},\qquad \|v\|^2=\sum_{x=1}^{N}x^{2k},
\end{equation}
the usual asymptotic estimate of the length of a shortest nonzero vector of $L$. (G6K uses the
exact volume of the unit ball, which gives a value larger by $3$ to $5\%$ in these dimensions;
the radius of Section~\ref{sec:sieve} and the ratios to $\GH$ in Section~\ref{sec:dbstats} use
this exact form.) For an integer vector
$c\in\Z^N$,
\begin{equation}\label{eq:norms}
\|c\|_2\le\|c\|_1\le\min\bigl(\|c\|_2^2,\ \sqrt N\,\|c\|_2\bigr).
\end{equation}
These inequalities are elementary. The second bound $\|c\|_1\le\|c\|_2^2$ holds because
$|c_x|\le c_x^2$, with equality exactly when every nonzero coefficient is $\pm1$; such a
relation has $T\le N$ terms but need not be balanced. The Lander--Parkin identity, for
instance, has $p=4$, $q=1$ and $\|c\|_1=\|c\|_2^2=5$. By \eqref{eq:norms}, a relation with
$T$ terms has $\|c\|_2^2\ge T^2/N$, and a relation of length close to $\GH(L)$ has at most
$\min(\GH^2,\sqrt N\,\GH)$ terms. Table~\ref{tab:gh} gives $\GH(L_{k,N})^2$; it is smallest at
$N=166$ for $k=20$ and $N=280$ for $k=30$, and to leading order in $N$ the minimum lies at
$N\approx2k(\ln N-1)$.

The Gaussian heuristic predicts Euclidean lengths and says nothing directly about $p$ and $q$.
It is an empirical observation, not a consequence of \eqref{eq:norms}, that the short vectors
of $L$ are nearly balanced. In the databases of Section~\ref{sec:dbstats}
(Table~\ref{tab:dbstats}) the median of $\min(p,q)$ is close to $T/2$, and at most $53$ of the
$5\times10^5$ to $1.2\times10^6$ vectors in each database have $\min(p,q)\le16$. The records
with small $m$ are much longer: for the 50 improved entries with $m\le4$ and $17\le k\le30$,
$\|c\|_2^2$ is between $1.3$ and $4.7$ times $\GH(L_{k,N})^2$, with $N$ the largest base. In
our runs they were found at small $N$; for example, our identities of types $(20,3,56)$ and
$(24,3,85)$ have largest bases $58$ and $59$. Table~\ref{tab:gh} therefore describes the
balanced relations that the sieve finds, not the unbalanced relations that make up most
records.
Section~\ref{sec:nsweep} compares values of $N$ at a fixed budget.

\begin{table}[htbp]
\centering\small
% generated by tools/make_tables.py -- do not edit
\begin{tabular}{rrrrrrrr}
\toprule
$k$ & $N=50$ & $N=70$ & $N=100$ & $N=150$ & $N=200$ & $N=250$ & best $N$ \\
\midrule
20 & 71 & 48 & 38 & 34 & 34 & 36 & 166 \\
22 & 97 & 61 & 45 & 39 & 38 & 39 & 188 \\
24 & 134 & 78 & 55 & 44 & 42 & 43 & 211 \\
26 & 184 & 100 & 66 & 51 & 47 & 46 & 233 \\
28 & 253 & 128 & 79 & 58 & 52 & 51 & 257 \\
30 & 348 & 164 & 95 & 66 & 58 & 55 & 280 \\
32 & 478 & 209 & 115 & 76 & 64 & 61 & 304 \\
\bottomrule
\end{tabular}

\caption{The heuristic $\GH(L_{k,N})^2$ of \eqref{eq:gh} and the value of $N$ that minimizes
it. By \eqref{eq:norms} it bounds the number of terms of a relation of length near the
shortest vector.}\label{tab:gh}
\end{table}

\subsection{The sieve database}\label{sec:sieve}
After BKZ-20 preprocessing we sieve the full lattice with G6K, progressively: we sieve the last
50 coordinates of the projected lattice, then extend the sieving context to the left one
dimension at a time and sieve again, until the full lattice is reached. A few early runs
started the full-dimensional sieve from random samples instead, which gives databases of the
same size and similar quality, more slowly (Section~\ref{sec:dbstats}).

We used two of G6K's sieves for the full lattice: \texttt{bgj1}, its variant of the sieve of
Becker, Gama and Joux \cite{BeckerGamaJoux2015}, and \texttt{bdgl}, the sieve of Becker, Ducas,
Gama and Laarhoven \cite{BeckerDucasGamaLaarhoven2016}. The GPU version of G6K, G6K-GPU
\cite{DucasStevensVanWoerden2021}, uses bucketing of the BDGL type, or in one run of the BGJ1
type, on the GPU. A sieve call stops when the database contains a given fraction, the
saturation ratio, of the expected number $\tfrac12\rho^{d/2}$ (up to sign) of lattice vectors of
length at most $\sqrt\rho\,\GH(L)$, where $\rho$ is the saturation radius and $\GH$ is computed
from the exact volume of the ball. In the CPU runs we set $\rho$ between $4/3$, G6K's default,
and $1.7$, the saturation ratio to $0.5$ and the database size to $3.2\,\rho^{d/2}$. The GPU
runs used saturation ratio $0.375$ and database size $2.77\,\rho^{d/2}$, the settings of
\cite{DucasStevensVanWoerden2021}, mostly with $\rho=4/3$. The radius is a stopping
criterion, not a bound on the database: in the CPU runs the database holds about
$3.2\,\rho^{d/2}$ vectors, so after saturation roughly $0.25/3.2\approx8\%$ of them lie
inside the radius and the rest are longer. In our measurements $92\%$ of the vectors lie
outside the radius and the longest have $\|c\|^2$ between about $1.7$ and $1.9$ times $\GH^2$
(Table~\ref{tab:dbstats}). If G6K signals that it failed to reach the target, we continue with
the current database; this did not happen in any of the \SatSieves{} sieve calls of the
experiments of Section~\ref{sec:experiments}, where we recorded it. All database vectors are
kept as moves. For $N=60$ and $\rho=1.5$ the database has $501{,}006$ vectors, for $N=90$ and
$\rho=1.38$ it has $5{,}367{,}467$; the GPU databases have about $4.2\times10^6$ vectors for
$N=100$ and $1.8\times10^7$ for $N=110$, and the two runs with $N=120$ kept the $6\times10^7$
shortest vectors (Section~\ref{sec:cost}).

By \eqref{eq:norms}, a relation with $T$ terms cannot be in a database whose longest vector has
squared length less than $T^2/N$. This necessary condition is weak. Our identity of type
$(28,1,126)$ has $T=127$ and largest base $69$, so it lies in $L_{28,70}$, and
$T^2/70\approx230$ is just below the largest squared length, $234$, in the four databases we
measured for $L_{28,70}$ (two seeds, progressive and random start); but its squared length is
$455$. Relations of this kind are not in databases of this size, and the records with small $m$
have to be reached by other means.

\subsection{Descent}\label{sec:descent}
Relations with small $\min(p,q)$ are longer than the database vectors. We reach them by a
greedy descent that uses the database $D$ as its set of moves (Algorithm~\ref{alg:descent}).
For a weight $w\ge0$ let
\[
f_w(c)=w\,\min(p(c),q(c))+p(c)+q(c),
\]
the objective of \texttt{ResultLC}. From each start the descent repeatedly replaces the current
relation $c$ by the best $c\pm u$ with $u\in D$, as long as $f_w$ strictly decreases, for at
most $t_{\max}=500$ steps. Ties are broken by the smallest index of $u$ in $D$, with $+$
before $-$. Every candidate $c\pm u$ examined is itself a relation, whether or not it lowers
$f_w$, and is used to update the frontier $F(m)$, the least $n$ found for each $m$. The
frontier therefore reflects the whole neighborhood of each path, not only the local minima where
the paths end; Section~\ref{sec:moveset} measures what this adds. Unless stated otherwise, the
starts are the rows of the reduced basis and 300 database vectors chosen at random, and the
weights are $\Omega=\{0,1,2,4,8,16,1000\}$. With $w=0$ the descent minimizes $p+q$; larger $w$
favors relations with smaller $\min(p,q)$. Two other kinds of start are used for particular
powers: coset starts, relations in a chosen coset of a sublattice of $L$
(Section~\ref{sec:congruence}), and known-identity starts, the identities with $m\le40$ and bases
at most $N$ in Meyrignac's database and, for $k\le32$, the 2003 archive, each alone and plus or minus random
database vectors (Section~\ref{sec:notimproved}).

The database is stored with one sign per vector (first nonzero entry positive), with
duplicates and the zero vector removed, as 16-bit integers; extraction checks that no
coefficient overflows. The descent works with 32-bit coefficients and 64-bit term counts, and
every identity is checked in exact arithmetic before it is stored. Both signs are used as
moves. The frontier is kept for $m<96$, with one
relation attaining each value. The descents for different starts and weights are independent
and run in parallel; each keeps its own frontier, and the frontiers are merged at the end. The
step limit is never reached in practice: over the \StepsTasks{} descents of the experiments of
Section~\ref{sec:experiments} the median number of accepted steps was \StepsMedian{} and the largest
\StepsMax{}. The cost of a step is $2|D|$ evaluations of $f_w$ in dimension $N$.

\begin{algorithm}[ht]
\caption{Descent with the sieve database as move set}\label{alg:descent}
\begin{algorithmic}[1]
\Require starts $S\subset L$, database $D=(u_1,\dots,u_r)\subset L$, weights $\Omega$, step limit $t_{\max}$
\State $F(m)\gets\infty$ for all $m$; update $F$ with every $s\in S$
\ForAll{$s\in S$ and $w\in\Omega$}
  \State $c\gets s$
  \For{$i=1,\dots,t_{\max}$}
    \State $(c^\ast,f^\ast)\gets(\bot,f_w(c))$ \Comment{$\bot$: no improving move found yet}
    \For{$j=1,\dots,r$ and $\sigma=+1,-1$}
      \State $c'\gets c+\sigma u_j$; \textbf{if} $c'=0$ \textbf{then continue}
      \State $m'\gets\min(p(c'),q(c'))$, $n'\gets\max(p(c'),q(c'))$; $F(m')\gets\min(F(m'),n')$
      \If{$f_w(c')<f^\ast$} $(c^\ast,f^\ast)\gets(c',f_w(c'))$ \EndIf
    \EndFor
    \If{$c^\ast=\bot$} \textbf{break} \Comment{local minimum}\EndIf
    \State $c\gets c^\ast$
  \EndFor
\EndFor
\State \Return $F$, together with a relation attaining each value
\end{algorithmic}
\end{algorithm}

\subsection{Combination of runs}\label{sec:combine}
Different values of $N$ and $\rho$ and different random choices give different relations. For
each $k$ we embed in $\Z^{120}$ every verified identity with bases at most $120$, both the
identities stored by our runs and those in Meyrignac's database and the 2003 archive, and
combine them by Algorithm~\ref{alg:combine}. Algorithm~\ref{alg:descent} is run with the pool
as both the starts and the moves. Because a relation that only ties a frontier value may
combine with another to give an improvement, the pool keeps, for each $m\le40$, a set
$\mathcal T_m$ of up to $\tau$ distinct relations attaining the current value $F(m)$, and since
the descent breaks ties by position, each round is run on several orderings of the pool. The
procedure stops after a round in which none of the sets $\mathcal T_m$ changes. It terminates:
for each $m\le40$, once $F(m)$ is finite it is a positive integer that never increases, and
while $F(m)$ is unchanged the set $\mathcal T_m$ only grows, up to $\tau$ elements, so each
$\mathcal T_m$ changes only finitely often. We use $\tau=64$, the given ordering and three
random ones, and at most 50 rounds; no run needed more than nine. On an early pool for $k=20$
the frontier with $\tau=8$ was three terms worse in total than with $\tau=32$, $128$ or $512$,
which all gave the same frontier; on later pools for $k=20$, $28$ and $30$ every $\tau$ from $8$
to $256$ gave the same frontier. The runs store only identities
that tie or improve the table, so the initial pool has between about 200 and 2{,}400 elements per
$k$.

\begin{algorithm}[ht]
\caption{Combination of a pool of relations}\label{alg:combine}
\begin{algorithmic}[1]
\Require initial pool $\mathcal P_0\subset L$, tie limit $\tau$, number of orderings $\eta$, weights $\Omega$
\State $F(m)\gets\infty$, $\mathcal T_m\gets\emptyset$ for all $m$; \Call{Offer}{$c$} for every $c\in\mathcal P_0$
\Repeat
  \State $\mathcal P\gets\mathcal P_0\cup\bigcup_m\mathcal T_m$; $\mathit{changed}\gets$ \textbf{false}
  \For{$i=1,\dots,\eta$}
    \State $\mathcal Q\gets\mathcal P$ in the given order ($i=1$) or a random order ($i>1$)
    \State run Algorithm~\ref{alg:descent} with starts $\mathcal Q$ and moves $\mathcal Q$, and for each
    descent and each $m$ \Call{Offer}{} the relation attaining its frontier value
    \State $\mathit{changed}\gets\mathit{changed}$ \textbf{or} some \Call{Offer}{} returned \textbf{true}
  \EndFor
\Until{\textbf{not} $\mathit{changed}$}
\Statex
\Function{Offer}{$c$} \Comment{$m=\min(p(c),q(c))$, $n=\max(p(c),q(c))$; ignored if $m>40$}
  \State \textbf{if} $n<F(m)$ \textbf{then} $F(m)\gets n$, $\mathcal T_m\gets\emptyset$
  \State \textbf{if} $n=F(m)$, $|\mathcal T_m|<\tau$ and $\pm c\notin\mathcal T_m$ \textbf{then} add $c$ to $\mathcal T_m$; \Return \textbf{true}
  \State \Return \textbf{false}
\EndFunction
\end{algorithmic}
\end{algorithm}

\subsection{Variants that did not improve the results}\label{sec:variants}
Three variants did not contribute to the final results.
\begin{enumerate}
\item A random $50$-element subset of $\{1,\dots,64\}$ gives a lattice of the same rank as
$L_{k,50}$ but with a larger determinant, and hence a larger $\GH$. Up to 25 such runs for
each $k$ from 19 to 30 improved only the entries $(22,8)$ and $(26,1)$, and full ranges of
bases later improved both further.
\item Sieving only the projection onto the last $d-f$ coordinates and lifting the database by
Babai's nearest-plane algorithm \cite{Babai1986}, as in ``dimensions for free'' \cite{Ducas2018}, reduces the
cost of the sieve, but the lifted relations are too long. For $k=28$, $N=70$ and $f=15$ the
median of $\|c\|_2^2$ rose from 217 to 285, and the frontier was worse for every $m\le16$.
\item Adding multiples of $\ell$ larger than $N$ to the bases, for primes $\ell$ with $(\ell-1)\mid k$,
increases $\det L$ more than it relaxes the congruence condition of
Section~\ref{sec:history}. It did not help for $k\in\{16,18,22,28,30\}$.
\end{enumerate}

%======================================================================
\section{Implementation and verification}\label{sec:impl}
%======================================================================

The search is written in Python using fpylll~0.6.4 with fplll~5.5.0 \cite{fpylll,fplll}
for LLL and BKZ~2.0 with fplll's default pruning strategies, G6K (commit \texttt{c71e084})
\cite{AlbrechtEtAl2019,G6Krepo} for CPU sieving and G6K-GPU-Tensor (hereafter G6K-GPU; branch
\texttt{python3lwe}, commit \texttt{e0e17ad}) \cite{DucasStevensVanWoerden2021,G6KGPUrepo}
for GPU sieving. The descent is compiled with Numba~0.68 for CPUs and as a CUDA kernel
(\texttt{numba-cuda}~0.20, CUDA~12.8--12.9) for GPUs. In the GPU kernel each thread evaluates one move
against eight relations, and the best move for each relation and the frontier are reduced with
64-bit atomic minima. The kernel breaks ties exactly as the CPU code does, and on a test lattice
both return the same frontier. G6K required two small portability fixes on ARM processors;
these are included with the code.

We checked every sieve used (\texttt{gauss}, which G6K uses in small dimensions, \texttt{bgj1}
and \texttt{bdgl} on CPUs; bucketing of the BGJ1 and BDGL types on GPUs)
against unpruned enumeration on random lattices of dimensions 50 to 56, where each returned
the exact shortest vector. Every identity reported is verified by evaluating both sides in
exact integer arithmetic before it is counted. The tables and the numbers in
Section~\ref{sec:results} are produced from the result files by a script that performs this
check; it checked \NumChecked{} identities, and none failed. Meyrignac's database and the
2003 archive were verified in the same way. The random choices of each run are determined by a
recorded seed, but the multithreaded sieves are not deterministic, so repeated runs give
different databases.

The data archive separates verification from discovery. It contains every identity that
improves an entry, the identities that lower the least known total $m+n$
(Section~\ref{sec:results}), the reference value and kind of every entry with
$17\le k\le100$ and $m\le16$, and a program of about 130 lines that uses only the Python
standard library. The program checks each identity in exact integer arithmetic, checks that no
base occurs on both sides, that the term counts agree with the stated type and that the
identity is primitive, and compares $n$ with the reference value of the entry
(Section~\ref{sec:problem}), or $m+n$ with the least total of an identity for the same $k$ in
Meyrignac's database and the 2003 archive. It reports the improvements separately for the two
kinds of entry. No part of the search software is needed.

Development used an Apple M3 (8 cores). Production runs used CPU nodes of the National
Research Platform (NRP; 12 or 16 cores per job) \cite{Weitzel2025NRP}, CPU nodes of the Turing
cluster at the Titan Supercomputing Center, California State University, Fullerton (32 cores
per job), a Turing node with four NVIDIA H200 GPUs and 1\,TB of memory, and two Turing nodes
with six NVIDIA T4 GPUs each. One NVIDIA L40 was used for GPU validation. Of the 24 CPU
searches with bases $1,\dots,100$ and the \texttt{bdgl} sieve, five (for $k=16$, $19$, $20$, $22$
and $23$) ran for six to nine hours without finishing, where the others took between 10 and 75
minutes; their
databases had stopped growing, and they were rerun on T4 GPUs with G6K-GPU. We did not find the
cause. Every GPU was checked against exact enumeration and the CPU descent before use.

%======================================================================
\section{Experiments}\label{sec:experiments}
%======================================================================

\subsection{Combination of known identities}\label{sec:controls}
Table~\ref{tab:controls} tests whether the results could have been obtained by combining the
identities already known. We applied Algorithm~\ref{alg:combine} separately for each
$17\le k\le30$, and the table has six columns: the identities of Meyrignac's database and the 2003 archive with bases at most $120$, as they are ($H$) and after
Algorithm~\ref{alg:combine} ($H^*$); the same identities with all bases after
Algorithm~\ref{alg:combine} ($H^*_{\mathrm{all}}$, in $\Z^{N}$ with $N$ the largest base, up to
1,921); the identities stored by our runs before any combination, as they are ($R$) and after
Algorithm~\ref{alg:combine} ($R^*$); and both together after Algorithm~\ref{alg:combine}
($RH^*$). Algorithm~\ref{alg:combine} applied to the known identities improves no entry, with or
without the restriction on the bases. The identities produced by the runs, before any
combination, improve \CtrlRImp{} of the \NumImproved{} entries, and Algorithm~\ref{alg:combine},
with or without the historical identities, does not improve any of them further. (The other
entries were improved by known-identity starts (Section~\ref{sec:notimproved}), whose
identities are not in the pools of Table~\ref{tab:controls}.) Earlier in the project, combination did lower several entries, but in
every case a later single run did at least as well. Since Algorithm~\ref{alg:combine} is a greedy
search, this does not show that no combination of the known identities improves an entry. It
shows that the improvements found without the known identities are present in the identities
found by individual runs, and that the combination step serves only to merge runs.

\begin{table}[htbp]
\centering\small
% generated by tools/controls.py -- do not edit
\begin{tabular}{lrrrrrr}
\toprule
pool & $H$ & $H^{*}$ & $H_{\mathrm{all}}^{*}$ & $R$ & $R^{*}$ & $RH^{*}$ \\
\midrule
entries improved & 0 & 0 & 0 & 191 & 191 & 191 \\
terms saved & 0 & 0 & 0 & 3,173 & 3,173 & 3,173 \\
\bottomrule
\end{tabular}

\caption{Entries with $m\le16$ and $17\le k\le30$ improved by each pool, without combination
($H$, $R$) and after Algorithm~\ref{alg:combine} (starred). $H$: identities of Meyrignac's
database and the 2003 archive with bases at most $120$; $H_{\mathrm{all}}$: all of them; $R$:
identities stored by our runs; $RH$: both.}\label{tab:controls}
\end{table}

\subsection{Stages and baselines}\label{sec:ablation}
Table~\ref{tab:ablation} compares the stages of the method with two baselines that use
reduction only, on the lattices $L_{k,60}$ for $k\in\{\AblK\}$, each with \AblSeeds{} seeds, on CPU nodes with
16 cores. The methods are as follows.
\begin{description}
\item[A] The 2003 procedure (the port of \texttt{SumsBKZ}): BKZ-20 followed by Yura's step,
repeated on new random orderings of the bases in 16 parallel processes.
\item[B] The sieve database alone ($\rho=1.5$, \texttt{bgj1}): the best $n$ for each $m$ among
its vectors.
\item[C] Sieve and descent (Algorithm~\ref{alg:descent}).
\item[D] Three runs of C with different seeds on the same reduced basis, combined by
Algorithm~\ref{alg:combine}.
\item[E] The same descent as C with a move set from reduction only: BKZ-20 followed by BKZ~2.0
with block size 40 on new random orderings in 16 parallel processes, for half the time of
method~D; all
distinct reduced rows (between $9\times10^4$ and $1.9\times10^5$) are the moves, and the
starts are the BKZ-20 rows and 300 random rows, as in C.
\end{description}
Methods A and E receive the wall time of method~D, the most expensive method; E spends half of
it on reduction and then runs the descent, which takes about a tenth of that time, so that E
uses about $0.6$ of the time of D in all, and C about a third. Since the methods reach different sets of entries, we compare
them on a common set: the entries that A, C, D and E reached in every seed. The column
``excess'' gives the mean of $n-n_{\mathrm{table}}$ over this set, averaged over the seeds,
with the standard deviation over the seeds.

The 2003 procedure (A) carried out between about 3,000 and 5,500 BKZ reductions per power and
seed. It improved at most three entries for $k=20$ and at most one for any other power. Its distance from the table grows with $k$:
its mean excess on the common entries is $1.9$ terms for $k=20$ and $105.5$ for $k=30$. The
database alone (B) improves at most one entry and reaches few. Adding the descent (C) gives
between 5 and 13 improvements per power for $k\in\{22,26,28,30\}$ in about a third of the time
of D, and combining three runs (D) adds up to four more, one or two on average. The
reduction-only move set (E) is the stronger baseline. For $k=24$ it is better than C and D, with 7 or 8
improvements against at most 3, and for $k=20$ it is better than C and comparable with D. For $k\in\{22,26,28,30\}$ it is
clearly worse than C, although C used less time. The spread over the seeds is smaller than
these differences. The descent with a large move set is what finds the records; the 2003
procedure, whose moves are the rows of a single reduced basis, is far behind except for $k=20$.
In this
comparison the sieve database is a better move set than the pool of BKZ-40 rows for four of the
six powers, but the two move sets differ in size, and three seeds per power make the comparison
indicative only; Section~\ref{sec:moveset} compares move sets of equal size.

We repeated the comparison on $L_{k,70}$ ($\rho=1.45$) for $k=24$ and $k=28$
(Table~\ref{tab:ablation70}). For $k=28$ the difference between the sieve and the
reduction-only move set is larger than at $N=60$. C improves 14 to 16 entries and E 7 or 8,
and the mean excess on the common entries is $-13.2$ for C against $-2.7$ for E. For $k=24$,
E improves more entries than C or D (4 to 6, against 2 and 2 to 5), but its mean excess on the
common entries is larger ($3.8$, against $2.5$ for C and $0.3$ for D). The power $k=24$ is the
one case in our experiments where the reduction-only move set is competitive at both
dimensions.


\begin{table}[htbp]
\centering\small
% generated by tools/ablation_table.py -- do not edit
\begin{tabular}{rlrrrrr}
\toprule
$k$ & method & time (s) & improved & reached & excess & median \\
\midrule
20 (6) & A & 160 & 0--3 & 12.7 & $1.9\pm0.5$ & $0.5$ \\
 & B & 11 & 0 & 1.0 & -- & -- \\
 & C & 49 & 0--2 & 10.3 & $3.1\pm1.3$ & $1.7$ \\
 & D & 158 & 1--5 & 16.0 & $1.2\pm0.3$ & $0.2$ \\
 & E & 94 & 3--5 & 16.0 & $1.3\pm0.2$ & $-0.2$ \\
\addlinespace
22 (14) & A & 169 & 0 & 14.3 & $31.1\pm3.5$ & $18.2$ \\
 & B & 12 & 0--1 & 5.3 & -- & -- \\
 & C & 60 & 5--9 & 16.0 & $1.1\pm0.4$ & $-0.2$ \\
 & D & 167 & 8--9 & 16.0 & $-0.3\pm0.5$ & $-0.5$ \\
 & E & 98 & 1--5 & 16.0 & $2.9\pm0.9$ & $1.3$ \\
\addlinespace
24 (8) & A & 168 & 0--1 & 8.3 & $24.7\pm0.9$ & $10.5$ \\
 & B & 11 & 0 & 0.0 & -- & -- \\
 & C & 58 & 0--2 & 12.7 & $8.3\pm1.7$ & $9.3$ \\
 & D & 167 & 0--3 & 16.0 & $4.8\pm1.4$ & $4.3$ \\
 & E & 99 & 7--8 & 16.0 & $1.4\pm0.8$ & $1.2$ \\
\addlinespace
26 (10) & A & 179 & 0--1 & 11.7 & $61.3\pm11.5$ & $47.3$ \\
 & B & 10 & 0 & 0.3 & -- & -- \\
 & C & 67 & 9--10 & 16.0 & $-0.9\pm1.9$ & $-0.7$ \\
 & D & 178 & 10--11 & 16.0 & $-3.3\pm1.4$ & $-2.7$ \\
 & E & 103 & 7 & 16.0 & $1.5\pm0.5$ & $1.7$ \\
\addlinespace
28 (13) & A & 140 & 0 & 13.3 & $84.7\pm2.7$ & $79.3$ \\
 & B & 9 & 0--1 & 2.0 & -- & -- \\
 & C & 48 & 10--12 & 16.0 & $-3.4\pm0.9$ & $-1.3$ \\
 & D & 139 & 12--13 & 16.0 & $-4.3\pm0.1$ & $-2.0$ \\
 & E & 80 & 5--8 & 16.0 & $3.4\pm1.7$ & $0.3$ \\
\addlinespace
30 (12) & A & 163 & 0 & 13.0 & $105.5\pm2.7$ & $93.0$ \\
 & B & 12 & 0 & 0.3 & -- & -- \\
 & C & 58 & 12--13 & 16.0 & $-2.1\pm0.3$ & $-3.3$ \\
 & D & 162 & 13--14 & 16.0 & $-4.2\pm1.3$ & $-4.0$ \\
 & E & 92 & 6--9 & 16.0 & $3.2\pm2.4$ & $-1.0$ \\
\bottomrule
\end{tabular}

\caption{The stages of the method and two baselines on $L_{k,60}$, \AblSeeds{} seeds each
(Section~\ref{sec:ablation}). Time is the mean wall time on 16 cores. ``Improved'' gives the
range over the seeds of the number of entries with $m\le16$ below the table; ``reached'' is
the mean number of entries for which an identity was found. ``Excess'' is the mean of
$n-n_{\mathrm{table}}$ over the entries reached by methods A, C, D and E in every seed, averaged
over the seeds, $\pm$ its standard deviation over the seeds; ``median'' is the median of
$n-n_{\mathrm{table}}$ over the same entries, averaged over the seeds. The number in parentheses
after $k$ is the number of these entries. A negative value means that the method is on average
below the table.}\label{tab:ablation}
\end{table}
\begin{table}[htbp]
\centering\small
% generated by tools/ablation_table.py -- do not edit
\begin{tabular}{rlrrrrr}
\toprule
$k$ & method & time (s) & improved & reached & excess & median \\
\midrule
24 (5) & A & 445 & 0 & 6.7 & $13.8\pm5.2$ & $11.7$ \\
 & B & 27 & 1 & 1.3 & -- & -- \\
 & C & 151 & 2 & 12.7 & $2.5\pm1.3$ & $4.3$ \\
 & D & 441 & 2--5 & 15.3 & $0.3\pm1.5$ & $-1.0$ \\
 & E & 260 & 4--6 & 16.0 & $3.8\pm0.7$ & $3.7$ \\
\addlinespace
28 (10) & A & 380 & 0 & 11.0 & $110.7\pm12.4$ & $98.0$ \\
 & B & 26 & 2--3 & 3.7 & -- & -- \\
 & C & 127 & 14--16 & 16.0 & $-13.2\pm1.1$ & $-8.7$ \\
 & D & 378 & 15--16 & 16.0 & $-14.3\pm0.2$ & $-12.3$ \\
 & E & 217 & 7--8 & 16.0 & $-2.7\pm1.7$ & $-0.2$ \\
\bottomrule
\end{tabular}

\caption{As Table~\ref{tab:ablation}, on $L_{k,70}$ with $\rho=1.45$.}\label{tab:ablation70}
\end{table}

\subsection{Move sets of equal size}\label{sec:moveset}
The sieve database and the move set of method~E differ in size as well as in content. To
separate the two we repeated the descent of method~C on $L_{k,60}$ for $k=20$, $24$ and $28$,
three seeds each, with the same starts and different move sets: the whole database, a random
sample of it and its shortest vectors, both of the size of method~E's move set for the same $k$
and seed, and random samples of $1/16$ to $1/2$ of it (Table~\ref{tab:moveset}). At equal size neither
move set is uniformly better: for $k=20$ and $24$ the rows of~E are clearly better (mean excess
$1.3$ and $1.5$, against $5.7$ to $14.2$ for the samples), while for $k=28$ a random sample is as
good as~E and the shortest sample is better. For $k=28$ the frontier
improves steadily with the size of the move set, from 2 to 4 improved entries with $1/16$ of
the database to 9 to 11 with all of it. What distinguishes the sieve is therefore mostly the
number of moves it supplies: on 16 cores it produced $5\times10^5$ vectors in about 11\,s,
while method~E needed about 80\,s of BKZ-40 reductions for $1.5\times10^5$ rows. (The starts of
method~E in Table~\ref{tab:ablation} were rows of~E rather than database vectors.)

The same experiment measures the frontier update of Section~\ref{sec:descent}. Recording only
the relations on the accepted path, or only the end points of the descents, raises the mean
excess from $1.8$ to $2.5$ and $3.0$ for $k=20$ and from $10.2$ to $13.3$ for $k=24$, and changes
little for $k=28$: the neighbors of the path supply the entries that are hardest to reach.

\begin{table}[htbp]
\centering\small
% generated by tools/moveset_table.py -- do not edit
\begin{tabular}{llrrrr}
\toprule
$k$ & move set & moves & improved & reached & excess \\
\midrule
20 (5) & sieve database & 501,006 & 0--1 & 11.0 & $1.8\pm0.6$ \\
 & random sample, $|E|$ & 156,074 & 0 & 10.3 & $5.7\pm1.4$ \\
 & shortest $|E|$ & 156,074 & 0 & 7.0 & $6.5\pm2.9$ \\
 & BKZ-40 rows (E) & 156,074 & 3--5 & 16.0 & $1.3\pm0.2$ \\
 & random $1/2$ & 250,503 & 0 & 12.3 & -- \\
 & random $1/4$ & 125,251 & 0 & 9.0 & -- \\
 & random $1/8$ & 62,625 & 0 & 8.3 & -- \\
 & random $1/16$ & 31,312 & 0 & 2.0 & -- \\
 & database, path only & 501,006 & 0--1 & 11.0 & $2.5\pm0.4$ \\
 & database, end points only & 501,006 & 0--1 & 10.3 & $3.0\pm0.5$ \\
\addlinespace
24 (6) & sieve database & 501,006 & 0 & 12.3 & $10.2\pm0.6$ \\
 & random sample, $|E|$ & 148,377 & 0--1 & 10.3 & $14.2\pm1.7$ \\
 & shortest $|E|$ & 148,377 & 0 & 13.0 & $9.4\pm2.7$ \\
 & BKZ-40 rows (E) & 148,377 & 7--8 & 16.0 & $1.5\pm0.7$ \\
 & random $1/2$ & 250,503 & 0--1 & 12.3 & -- \\
 & random $1/4$ & 125,251 & 0 & 11.7 & -- \\
 & random $1/8$ & 62,625 & 0 & 12.0 & -- \\
 & random $1/16$ & 31,312 & 0 & 8.0 & -- \\
 & database, path only & 501,006 & 0 & 10.3 & $13.3\pm1.8$ \\
 & database, end points only & 501,006 & 0 & 8.7 & $13.3\pm1.8$ \\
\addlinespace
28 (16) & sieve database & 501,006 & 9--11 & 16.0 & $-5.3\pm1.0$ \\
 & random sample, $|E|$ & 131,811 & 6--10 & 16.0 & $2.5\pm3.4$ \\
 & shortest $|E|$ & 131,811 & 8--10 & 16.0 & $-1.9\pm2.0$ \\
 & BKZ-40 rows (E) & 131,811 & 5--8 & 16.0 & $2.5\pm1.3$ \\
 & random $1/2$ & 250,503 & 7--10 & 16.0 & -- \\
 & random $1/4$ & 125,251 & 6--10 & 16.0 & -- \\
 & random $1/8$ & 62,625 & 4--6 & 16.0 & -- \\
 & random $1/16$ & 31,312 & 2--4 & 16.0 & -- \\
 & database, path only & 501,006 & 9--10 & 16.0 & $-5.2\pm1.0$ \\
 & database, end points only & 501,006 & 9--10 & 16.0 & $-5.1\pm1.0$ \\
\bottomrule
\end{tabular}

\caption{The descent of method~C on $L_{k,60}$, three seeds per power, with different move sets
and the same starts (the BKZ-20 rows and 300 database vectors). ``Moves'': mean size of the move
set, where $|E|$ is the size of method~E's move set for the same $k$ and seed. ``Improved'' and
``reached'' as in Table~\ref{tab:ablation}. ``Excess'': mean of $n-n_{\mathrm{table}}$ over the
entries reached by all six variants with a value in this column, with the standard deviation
over the seeds; the number in parentheses after $k$ is the number of these entries. The last two
rows update the frontier only on the accepted path or only at the end points of each descent.
The row for~E is the run of method~E in Table~\ref{tab:ablation}, whose starts were rows of~E;
its excess is taken over the entries of this table and so differs slightly from
Table~\ref{tab:ablation}.}\label{tab:moveset}
\end{table}

\subsection{Properties of the sieve database}\label{sec:dbstats}
Table~\ref{tab:dbstats} describes one database for each lattice of Section~\ref{sec:ablation}
and for $L_{20,70}$. The fraction outside the saturation radius and the
achieved saturation are fixed by G6K's parameters, as explained in Section~\ref{sec:sieve}.
The databases are balanced in the sense of Section~\ref{sec:gh}: the median of $\min(p,q)$ is
about $T/2$, and the least $\min(p,q)$ in a database of $5\times10^5$ to $1.2\times10^6$ vectors is
between 9 and 19. The database alone therefore reaches few entries of the table (method B).

The progressive sieve was about 2.4 to 4 times faster than the full-dimensional sieve from
random samples on the six instances of Table~\ref{tab:progressive}, and its database was
slightly longer (median $\|c\|^2$ larger by $0.4$ to $1.2\%$). The frontiers obtained by the same descent
from the two databases differ by amounts in both directions, without a consistent sign, which
is what one expects from two random databases of the same size and similar lengths.

\begin{table}[htbp]
\centering\footnotesize
\setlength{\tabcolsep}{4pt}
% generated by tools/ablation_table.py -- do not edit
\begin{tabular}{rrrrrrrrr}
\toprule
$k$ & $N$ & $\rho$ & size & out & $\frac{\max\|c\|^2}{\GH^2}$ & $T$ med., max & $\min(p,q)$ min, med. & sat.\\
\midrule
20 & 60 & 1.5 & 501,006 & 92\% & 1.83 & 60, 73 & 14, 29 & 0.51 \\
22 & 60 & 1.5 & 501,006 & 92\% & 1.88 & 69, 85 & 9, 34 & 0.50 \\
24 & 60 & 1.5 & 501,006 & 92\% & 1.81 & 80, 95 & 19, 37 & 0.50 \\
26 & 60 & 1.5 & 501,006 & 92\% & 1.87 & 93, 111 & 16, 42 & 0.50 \\
28 & 60 & 1.5 & 501,006 & 92\% & 1.84 & 107, 127 & 15, 50 & 0.50 \\
30 & 60 & 1.5 & 501,006 & 92\% & 1.86 & 124, 150 & 14, 58 & 0.50 \\
20 & 70 & 1.45 & 1,181,262 & 92\% & 1.79 & 57, 71 & 18, 28 & 0.50 \\
24 & 70 & 1.45 & 1,181,262 & 92\% & 1.70 & 74, 91 & 16, 34 & 0.50 \\
28 & 70 & 1.45 & 1,181,262 & 92\% & 1.69 & 96, 114 & 12, 47 & 0.50 \\
\bottomrule
\end{tabular}

\caption{The sieve databases, one per lattice (the first seed). ``Out'': fraction of vectors
with $\|c\|^2>\rho\,\GH^2$; ``sat.'': number of vectors within the radius divided by
$\tfrac12\rho^{d/2}$ (G6K's target is $0.5$). $\GH$ uses the exact ball volume.}\label{tab:dbstats}
\end{table}

\begin{table}[htbp]
\centering\small
% generated by tools/ablation_table.py -- do not edit
\begin{tabular}{rrrrrrr}
\toprule
& & \multicolumn{2}{c}{sieve time (s)} & \multicolumn{2}{c}{median $\|c\|^2/\GH^2$} & \\
\cmidrule(lr){3-4}\cmidrule(lr){5-6}
$k$ & seed & progr. & random & progr. & random & $\sum_m (n_{\mathrm P}-n_{\mathrm R})$ \\
\midrule
20 & 1 & 24 & 83 & 1.619 & 1.600 & $+11$ \\
20 & 2 & 26 & 61 & 1.619 & 1.600 & $+19$ \\
24 & 1 & 32 & 118 & 1.581 & 1.569 & $+8$ \\
24 & 2 & 27 & 92 & 1.581 & 1.569 & $+16$ \\
28 & 1 & 31 & 78 & 1.574 & 1.567 & $-7$ \\
28 & 2 & 32 & 131 & 1.574 & 1.567 & $-24$ \\
\bottomrule
\end{tabular}

\caption{Progressive sieving against a full-dimensional sieve from random samples on
$L_{k,70}$, $\rho=1.45$, 16 cores. The last column is the total difference of the frontiers
(progressive minus random) over the entries $m\le16$ reached by both; it varies in sign.}
\label{tab:progressive}
\end{table}

\subsection{The choice of \texorpdfstring{$N$}{N}}\label{sec:nsweep}
Table~\ref{tab:gh} predicts that the shortest relations, and hence the shortest balanced
identities, need $N$ between about 160 and 280 for $20\le k\le30$. The records with small $m$
behave differently. To
compare values of $N$ we gave each pair $(k,N)$ with $k\in\{20,24,28\}$ and
$N\in\{50,60,70,80,90\}$ the same budget of one hour on 16 cores, repeating method~C on new
random orderings of the bases while time remained. The saturation radius was $1.5$ for
$N\le60$ and decreased to $1.38$ at $N=90$, so the comparison also varies $\rho$.
Table~\ref{tab:nsweep} gives the result.
The best $N$ grows with $k$: it is $60$ for $k=20$, where $N=60$ improves 11 entries, $70$ to $80$
for $k=24$, and $80$ for $k=28$, where $N=80$ is best or nearly best for every $m$ and $N=50$
is far behind. For the smaller powers the entries with small $m$ favor small $N$. For
$k=20$ only $N=50$ and $60$ reach any entry with $m\le4$ within the hour, and for $k=24$ only
$N=50$ reaches $m=1$. Finally, $N=90$ completes only three or four runs in the budget and is
worse than $N=80$ for every power. All of these values lie far below the dimensions at which
$\GH$ is smallest (Table~\ref{tab:gh}). For the unbalanced relations that make up the
table, the useful dimension is set by the cost of a run and by the reach of the descent, not
by the length of the shortest vectors. The experiment measures the combined effect of the two,
since a larger $N$ allows fewer runs within the budget; it does not separate them.

\begin{table}[htbp]
\centering\footnotesize
\setlength{\tabcolsep}{4pt}
% generated by tools/nsweep_table.py -- do not edit
\begin{tabular}{rrrrrrrr}
\toprule
$k$ & $N$ & $\rho$ & runs & improved & $m\le4$ reached & gap, $m\le4$ & gap, $5\le m\le16$ \\
\midrule
20 & 50 & 1.5 & 449 & 4 & 4 & -- & 47 \\
 & 60 & 1.5 & 104 & 11 & 4 & -- & 25 \\
 & 70 & 1.45 & 26 & 3 & 0 & -- & 25 \\
 & 80 & 1.4 & 14 & 3 & 0 & -- & 25 \\
 & 90 & 1.38 & 4 & 1 & 0 & -- & 58 \\
\addlinespace
24 & 50 & 1.5 & 613 & 5 & 4 & 11 & 114 \\
 & 60 & 1.5 & 60 & 9 & 3 & 15 & 59 \\
 & 70 & 1.45 & 20 & 10 & 3 & 19 & 53 \\
 & 80 & 1.4 & 16 & 9 & 3 & 13 & 47 \\
 & 90 & 1.38 & 4 & 4 & 2 & 30 & 98 \\
\addlinespace
28 & 50 & 1.5 & 642 & 6 & 4 & 115 & 357 \\
 & 60 & 1.5 & 74 & 14 & 4 & 50 & 179 \\
 & 70 & 1.45 & 23 & 16 & 4 & 26 & 32 \\
 & 80 & 1.4 & 13 & 16 & 4 & 7 & 7 \\
 & 90 & 1.38 & 3 & 16 & 4 & 9 & 61 \\
\bottomrule
\end{tabular}

\caption{Method C repeated for one hour on 16 cores for each $k$ and $N$. ``Runs'' is the
number of runs completed; ``improved'' counts entries $m\le16$ below the table, and ``$m\le4$
reached'' the entries with $m\le4$ for which an identity was found. The two gap columns give the
distance to the best $N$, $\sum_m\bigl(n_N(m)-\min_{N'}n_{N'}(m)\bigr)$ over the entries reached
at every $N$, separately for $m\le4$ and $5\le m\le16$.}\label{tab:nsweep}
\end{table}

\subsection{Saturation radius}
For $k=20$ and $N=50$, raising $\rho$ from $4/3$ to $1.5$ and $1.7$ increases the database
from $3{,}681$ to $65{,}976$ and $1{,}416{,}279$ vectors and improves the frontier. With
$\rho=1.5$ the search took 7\,s on the M3 and found $(20,3,60)$, one term below the table entry; with
$\rho=1.7$ it took 150\,s and also found $(20,2,61)$, below the table entry $62$, and matched the
table for $m=4$ and $m=5$.

\subsection{Cost}\label{sec:cost}
The heuristic cost of a sieve grows exponentially in the dimension, as $2^{0.292d+o(d)}$ for
the BDGL sieve \cite{BeckerDucasGamaLaarhoven2016}, and with a larger exponent for the
\texttt{bgj1} sieve of G6K \cite{AlbrechtEtAl2019,BeckerGamaJoux2015}, used in most of our CPU
runs. The descent costs $2|D|$ evaluations per
step for each of the $(N+299)|\Omega|$ descents, and with $|D|$ growing like $\rho^{d/2}$ it is the
larger cost on CPUs for $N\le70$; from $N=80$ on, the sieve costs more. On 32 CPU cores a
search with $N=100$ ($\rho=4/3$, \texttt{bdgl}) took about an hour per power, most of it in the sieve. On
GPUs both stages are much faster, and the descent becomes a small part of the cost: for $k=28$
and $N=70$ it took 6\,s on an L40, against 207\,s on the eight cores of the M3. For $N=120$ the
database at $\rho=4/3$ has about $7.5\times10^7$ vectors. Two such runs, for $k=28$ and $30$ on
the node with four H200 GPUs, kept the $6\times10^7$ shortest of them as moves. They lowered our
previous best values for $(28,15)$ from $70$ to $48$, $(28,16)$ from $52$ to $46$, $(30,12)$ from
$102$ to $82$ and $(30,16)$ from $76$ to $50$, below every other run, and lowered no other entry
with $m\le16$. For $N=140$ the database would have about $10^9$ vectors, which would need roughly
the whole memory of the node.

All runs reported in this paper, including those of Section~\ref{sec:high}, used about
\CoreHoursTotal{} CPU core-hours (\CoreHoursNRP{} on the National Research Platform and
\CoreHoursTuring{} on Turing) and \GPUHours{} GPU-hours on the H200 and T4 GPUs, counted as
wall time multiplied by the cores or GPUs allocated. Development and the early runs on the 8-core M3
are not included.

%======================================================================
\section{Results}\label{sec:results}
%======================================================================

This section and the next report the best identities found by any of the searches, which
differ in $N$, $\rho$, method and number of repetitions (Table~\ref{tab:runs} lists those of this
section).
They do not measure what a single method yields under a fixed budget; that is the subject of
Section~\ref{sec:experiments}.

We count an entry as improved when we have a primitive identity of type $(k,m,n)$ with $n$
smaller than the reference value of the entry (Section~\ref{sec:problem}). Every value we report as an improvement is also below the bound $n(m')+(m-m')$
that the table's rule derives from each entry $(k,m')$ with $m'<m$, so no improvement merely
fills a gap that the table's own rule closes. We do not apply the table's rule $n\mapsto n+i$
for the entry $(k,m+i)$ to our own identities; doing so would improve \NumClosureExtra{} further entry, \ClosureExtra, and would
lower \NumClosureLower{} of the values reported below.

The searches are summarized in Table~\ref{tab:runs}. Of the \NumCells{} entries with
$m\le16$ and $17\le k\le30$, \NumImproved{} are improved and \NumTies{} are matched, with a
total saving of \TermsSaved{} terms. The improvements are of two kinds
(Table~\ref{tab:perk}). Of the \NumCellsIdent{} entries attained by an identity,
\NumImprovedIdent{} are improved, saving \TermsSavedIdent{} terms. Of the \NumCellsDerived{}
entries that are derived bounds, \NumDerived{} are improved, saving \TermsSavedDerived{} terms;
for these entries the database had no identity for this $(k,m)$, or only a longer one, and ours
is the first identity below the bound. Three of them, \NoIdentCells{}, have no identity of any length
in the database. \TiesDerivedSentence{} The largest saving is for \MaxGainCell, from $\MaxGainOld$ to $\MaxGainNew$.
Every entry is improved for $k\in\{\FullKList\}$, and all new identities use bases at most
\MaxBase. All of them are primitive, and none has a common factor in its multiplicities.
Table~\ref{tab:grid} in Appendix~\ref{app:grid} gives every entry; the identities are in the
data archive.

\begin{table}[htbp]
\centering\footnotesize
\setlength{\tabcolsep}{3.5pt}
\begin{tabular}{>{\raggedright\arraybackslash}p{2.5cm}l>{\raggedright\arraybackslash}p{2.6cm}>{\raggedright\arraybackslash}p{1.9cm}>{\raggedright\arraybackslash}p{3.4cm}}
\toprule
search & $N$ & $\rho$, sieve & hardware & powers (runs per power) \\
\midrule
exploratory searches\textsuperscript{a} & 50--90 & 4/3--1.7; \texttt{bgj1}, \texttt{bdgl} & M3, NRP & $17\le k\le30$ (1 to 25 per $N$); $k=31$, $32$ (1) \\
experiments of Section~\ref{sec:experiments} & 50--90 & 1.38--1.5; \texttt{bgj1} & NRP & $k=20$, $22$, \dots, $30$ \\
searches & 100 & 4/3; \texttt{bdgl}, G6K-GPU & Turing CPU, T4, H200 & $k=16$, $19\le k\le30$ (1; 2 for $k=24$, 5 for $k=28$) \\
searches & 120 & 4/3; G6K-GPU & H200 & $k=28$, $30$ (1) \\
repetitions\textsuperscript{b} & 60--110 & 4/3--1.5; G6K-GPU & T4, H200 & $17\le k\le32$ (up to 6) \\
coset and known-identity starts\textsuperscript{c} & 48--110 & 4/3--1.7; G6K-GPU & T4, H200 & $17\le k\le21$, $24\le k\le26$, $32$ (up to 10 per $N$ and $\rho$) \\
\bottomrule
\end{tabular}
\caption{Searches whose identities are counted in Section~\ref{sec:results}, all with bases
$1,\dots,N$ except as stated in (a). Except where noted, each run used BKZ-20 preprocessing (BKZ-30 in two runs on the M3), the whole database as
moves, the reduced basis and 300 database vectors as starts, and the weights $\Omega$; the GPU
runs used the sieve settings of Section~\ref{sec:sieve}, taken from
\cite{DucasStevensVanWoerden2021,G6KGPUrepo}, except one run with $N=100$ on the H200, which used
the defaults of G6K-GPU. (a) Also
random 50-element subsets of $1,\dots,64$ and bases with added multiples
(Section~\ref{sec:variants}). (b) Repetitions of each power's most productive search with $N$ and
$\rho$ as in its earlier best run, and then with $N=110$; $k=25$ and $k=31$ repeated the methods
of Section~\ref{sec:high}. (c) Coset starts and known-identity starts
(Sections~\ref{sec:notimproved} and~\ref{sec:congruence}), each repetition with a new database.
The identities were then combined as in Section~\ref{sec:combine}.}\label{tab:runs}
\end{table}

\begin{table}[htbp]
\centering\small
% generated by tools/make_tables.py -- do not edit
\begin{tabular}{rrrrrr}
\toprule
& \multicolumn{3}{c}{entries improved} & & \\
\cmidrule(lr){2-4}
$k$ & identity & bound & total & terms saved & largest saving \\
\midrule
17 & 7 & 2 & 9 & 21 & 6 \\
18 & 4 & 0 & 4 & 22 & 15 \\
19 & 10 & 3 & 13 & 49 & 6 \\
20 & 11 & 2 & 13 & 66 & 10 \\
21 & 14 & 1 & 15 & 130 & 16 \\
22 & 13 & 3 & 16 & 118 & 24 \\
23 & 15 & 1 & 16 & 240 & 23 \\
24 & 12 & 4 & 16 & 198 & 24 \\
25 & 13 & 3 & 16 & 185 & 26 \\
26 & 12 & 3 & 15 & 235 & 31 \\
27 & 9 & 7 & 16 & 419 & 45 \\
28 & 15 & 1 & 16 & 449 & 57 \\
29 & 10 & 6 & 16 & 559 & 50 \\
30 & 11 & 5 & 16 & 575 & 58 \\
\midrule
all & 156 & 41 & 197 & 3,266 & 58 \\
\bottomrule
\end{tabular}

\caption{Improvements by power, separated into entries attained by an identity in the table
and entries that are derived bounds, with the total number of terms saved and the largest
saving in a single entry.}\label{tab:perk}
\end{table}

Two of the shortest new identities show the typical size of the solutions. The first, of
type $(17,8,22)$, improves the table entry $24$:
\begin{multline*}
90^{17}+83^{17}+58^{17}+48^{17}+42^{17}+36^{17}+24^{17}+3^{17}\\
=88^{17}+84^{17}+82^{17}+75^{17}+68^{17}+62^{17}+55^{17}+51^{17}+41^{17}+40^{17}\\
+39^{17}+38^{17}+2\cdot27^{17}+26^{17}+23^{17}+21^{17}+3\cdot11^{17}+9^{17}+5^{17}.
\end{multline*}
The second, of type $(17,9,21)$, improves the entry $23$:
\begin{multline*}
85^{17}+77^{17}+73^{17}+67^{17}+65^{17}+64^{17}+54^{17}+21^{17}+6^{17}\\
=84^{17}+76^{17}+75^{17}+2\cdot72^{17}+71^{17}+70^{17}+62^{17}+60^{17}+2\cdot56^{17}\\
+50^{17}+42^{17}+41^{17}+27^{17}+24^{17}+2\cdot23^{17}+16^{17}+14^{17}+8^{17}.
\end{multline*}

\subsection{Variation between runs}\label{sec:variation}
The sieve, the random starts of the descent and the random choices of the completions make
repeated runs differ. After all the searches described in this paper, we repeated for every
power the method that had produced most of its best identities, in two rounds of two runs
each with new random choices, and then in two rounds with bases up to $110$, for 75 and then 74
of the 84 powers (Table~\ref{tab:runs} and Section~\ref{sec:high}). We say that a round lowers
an entry when it finds an identity shorter than the best of all earlier runs, and, as before,
that it improves an entry when it goes below the reference value for the first time. The first
round lowered more than 700 of the entries for $17\le k\le100$ and improved 13 more. The second
lowered about 430, saving more than 27,000 terms, and improved two more. The first round with
$N=110$ lowered about 970, saving about 272,500 terms, and improved nine more; for 35 of the 45
powers above $55$ it lowered every entry. The second round with $N=110$ lowered about 450,
saving about 25,000 terms, but improved no further entry. The values we report are the best over all runs; they are not
converged, and further runs would lower many of them.

\subsection{Entries not improved}\label{sec:notimproved}
Of the \NumNotImproved{} entries with $17\le k\le30$ that are not improved, three have record
identities with bases beyond the range of our searches or at its edge: $(17,1)$ and $(26,14)$,
with largest bases $208$ and $149$, and $(20,1)$, with $106$. Only the known-identity starts
with bases up to $110$ covered $(20,1)$, without improving it. All three records involve
Wr\'oblewski's congruence methods. Each of the others is a derived bound or has a record identity
with bases at most $96$, within the range searched, which our runs did not reach. Most of these
entries are for $k=18$, with 11, and $k=17$, with five; $k=19$, $20$, $21$ and $26$ have at most
three each. For $k=17$ and $19$ the first moment is forced in the sense of
Section~\ref{sec:congruence}, with $M_1(k)=510$ and $798$; for $k=18$ no moment is forced and we
have no explanation. Its sieve database differs from those of its neighbors in one statistic we
measured: on $L_{18,100}$ the largest quarter of the bases carries $16\%$ of the mass
$\sum|c_x|$, against $19$ to $22\%$ for $k=17$, $19$, $20$ and $21$.

Two methods of Section~\ref{sec:congruence} work much better at small $N$. The known-identity
starts and the coset starts were first run with the databases of $L_{k,100}$ and $L_{k,110}$
(and, for $k=50$, of $L_{50,117}$). With $N$ between $48$ and $96$, just above the largest bases
of the records, a saturation radius between $1.4$ and $1.7$, and up to 10 repetitions with new
databases for each $N$ and $\rho$, they improved 16 entries that no search with larger $N$ had
improved: $(17,2)$; $(18,3)$, $(18,4)$, $(18,5)$ and $(18,8)$, the first improvements for $k=18$;
$(19,5)$ and $(19,6)$; $(24,7)$ and $(24,11)$; $(25,3)$ to $(25,6)$; and $(32,1)$, $(32,2)$ and
$(32,15)$, with $(32,1)$ falling from $230$ to $194$. They also lowered 23 entries that were
already improved. With $\rho=4/3$ the same runs improved only $(25,6)$, and later repetitions
continued to find improvements, so these entries are not converged either. Every entry is now
improved for $k=24$, $25$ and $32$.

For $k=31$ the record identities have largest bases between $127$ and $265$, beyond the bases
of our other searches ($N\le100$); the completions of Section~\ref{sec:completion} improve all 16
entries. For $k=32$ the searches with $N\le110$ and the starts of the previous paragraph improve
all 16 entries, by \KThirtyTwoSaved{} terms in all; Table~\ref{tab:gh} suggests that
$N\le110$ is far from optimal for this power.

\subsection{Least total}\label{sec:leasttotal}
The table also records, for each $k$, the identities with the least total $m+n$, which may have
$m>16$. For $k=21$, $23\le k\le30$ and $k=32$, and for $33\le k\le38$, $43\le k\le48$, and
$k=51$, $53$ and $56$ in the second table, one of our identities has fewer terms in total than every identity
for the same $k$ in Meyrignac's database and the 2003 archive, and than the least total listed
in his table; for example $(28,27,31)$ against
$(28,31,32)$, and $(43,40,97)$ against $(43,59,106)$. For $k=17$, $19$, $20$, $22$ and $39$ we
match the least total. For the other powers the least totals belong to nearly balanced
identities, $(18,13,15)$ for $k=18$ and identities with $m\ge33$ for $k=31$ and $k\ge40$, which
our searches did not target. These identities are in the data
archive.

\subsection{Comparison with other sources}\label{sec:sources}
Our improvements are measured against Meyrignac's table and database. To check that they are
also new with respect to the literature, we searched MathSciNet for the years 2015--2026: all
items classified in 11D41 or 11D72, and title searches for equal sums, like powers, Euler's
conjecture, and the Prouhet--Tarry--Escott and multigrade problems, about 1,200 items in all.
We also searched Google Scholar from 2015 on (the phrase ``equal sums of like powers'', and
Meyrignac, his site, ``Euler's extended conjecture'', the conjecture of Lander, Parkin and
Selfridge, and lattice methods for like powers), arXiv, zbMATH Open, the OEIS and the main web
collections on equal sums of like powers
\cite{WroblewskiTables,PiezasCollection,ChenShuwenESLP,PerrenetPowersums}. None of these sources
reports an identity with $k\ge11$ that improves an entry of the table; the closest publications
treat small degrees, for example \cite{Choudhry2016,Wagstaff2022,Braun2026}.

%======================================================================
\section{The table for \texorpdfstring{$33\le k\le100$}{33 <= k <= 100}}\label{sec:high}
%======================================================================

Meyrignac's second table lists the least known $n$ for $33\le k\le100$ and $m\le16$. Of its
\HNumCells{} entries, \HNumCellsIdent{} are attained by an identity and \HNumCellsDerived{} are
derived bounds. Of the identities, 518 date from 2001 and 84 from 2002; most were found by
Torbj\"orn Alm with \texttt{SumsBKZ}, and the median of their largest bases is $115$.

\subsection{Searches and results}\label{sec:highsearch}
For these powers the bound \eqref{eq:norms} is governed by $\sqrt N\,\GH(L)$ rather than by
$\GH(L)^2$, since the coefficients of short relations are no longer small. At $N=100$ a
relation of length $\GH(L)$ has at most about $250$ terms for $k=50$ and $2{,}500$ for $k=100$,
while the table entries for these powers are between $700$ and $1{,}200$ and between
$11{,}000$ and $14{,}000$. We therefore ran one search with bases $1,\dots,100$ for each
$33\le k\le100$, which we call the plain search: method~C with $\rho=4/3$, on 32 CPU cores
(\texttt{bdgl}) for $33\le k\le41$ and on one NVIDIA T4 GPU with G6K-GPU for $42\le k\le100$. The repetitions of Section~\ref{sec:results} used
each power's most productive method, the last of them with bases up to $110$.

For the powers where this search failed we also searched with bases $1,\dots,60$ and
$1,\dots,70$, which improved every entry for $k=40$ and $42$, used coset starts and, for $k=50$,
known-identity starts (Section~\ref{sec:congruence}),
completed Wr\'oblewski's congruence solutions (Section~\ref{sec:completion}) and constructed
congruence solutions of our own (Section~\ref{sec:structured}). Together these improve \HNumImproved{} of the \HNumCells{}
entries, \HNumImprovedIdent{} of them attained by an identity and \HNumDerived{} derived
bounds, and every entry for \HNumFullK{} of the \HNumK{} powers (Table~\ref{tab:perkhigh}). Because many of the old
entries are far from what the lattice offers, the savings are large: \HTermsSaved{} terms in
all, and \HMaxGain{} for the entry \HMaxGainCell, from \HMaxGainOld{} to \HMaxGainNew. \HNoneSentence{}
Table~\ref{tab:contrib} separates the methods. Most improvements come from the sieve search
itself, with all values of $N$ and all repetitions. Of the other methods, the coset starts and
the completions of Wr\'oblewski's solutions each improve more than 90 entries that no other
method improves, the known-identity starts 14 entries for $k=50$, and the completions of
structured solutions the 31 improved entries for $k=75$, $85$ and $97$.

\begin{table}[htbp]
\centering\scriptsize
\setlength{\tabcolsep}{3pt}
% generated by tools/make_tables.py -- do not edit
\begin{tabular}{rrrrr@{\hspace{1.5em}}rrrrr}
\toprule
$k$ & ident. & bound & saved & largest & $k$ & ident. & bound & saved & largest \\
\midrule
33 & 9 & 7 & 1,249 & 124 & 67 & 9 & 7 & 24,741 & 2,377 \\
34 & 14 & 2 & 1,609 & 114 & 68 & 10 & 6 & 27,337 & 3,132 \\
35 & 12 & 4 & 2,197 & 188 & 69 & 11 & 5 & 32,673 & 4,033 \\
36 & 13 & 3 & 1,936 & 177 & 70 & 10 & 6 & 26,772 & 1,869 \\
37 & 11 & 5 & 933 & 87 & 71 & 8 & 8 & 33,155 & 3,450 \\
38 & 8 & 8 & 2,547 & 244 & 72 & 8 & 8 & 29,238 & 2,274 \\
39 & 9 & 7 & 2,975 & 340 & 73 & 5 & 11 & 19,757 & 1,285 \\
40 & 10 & 6 & 2,253 & 352 & 74 & 6 & 10 & 22,443 & 1,899 \\
41 & 9 & 7 & 2,950 & 303 & 75 & 5 & 7 & 52,270 & 4,558 \\
42 & 10 & 6 & 1,733 & 296 & 76 & 8 & 8 & 42,653 & 4,850 \\
43 & 13 & 3 & 4,248 & 293 & 77 & 13 & 3 & 35,249 & 3,252 \\
44 & 11 & 5 & 4,931 & 447 & 78 & 7 & 9 & 40,971 & 3,554 \\
45 & 12 & 4 & 5,815 & 544 & 79 & 10 & 6 & 45,583 & 3,627 \\
46 & 12 & 4 & 4,897 & 479 & 80 & 9 & 7 & 46,658 & 5,591 \\
47 & 12 & 4 & 6,426 & 527 & 81 & 7 & 9 & 29,848 & 3,412 \\
48 & 8 & 8 & 6,953 & 608 & 82 & 9 & 7 & 67,814 & 7,001 \\
49 & 8 & 8 & 6,092 & 515 & 83 & 6 & 10 & 52,746 & 5,432 \\
50 & 14 & 2 & 6,916 & 587 & 84 & 10 & 6 & 51,807 & 5,081 \\
51 & 12 & 4 & 8,491 & 653 & 85 & 2 & 7 & 20,037 & 2,647 \\
52 & 9 & 7 & 12,086 & 1,292 & 86 & 4 & 12 & 56,921 & 3,918 \\
53 & 9 & 7 & 10,068 & 891 & 87 & 5 & 11 & 54,533 & 6,133 \\
54 & 12 & 4 & 12,336 & 1,178 & 88 & 11 & 5 & 68,457 & 9,654 \\
55 & 8 & 8 & 9,601 & 708 & 89 & 7 & 9 & 61,912 & 7,786 \\
56 & 11 & 5 & 12,614 & 1,107 & 90 & 9 & 7 & 76,021 & 7,681 \\
57 & 13 & 3 & 15,320 & 1,801 & 91 & 9 & 7 & 67,427 & 7,456 \\
58 & 11 & 5 & 15,191 & 1,421 & 92 & 7 & 9 & 82,239 & 6,809 \\
59 & 7 & 9 & 15,368 & 1,419 & 93 & 9 & 7 & 80,990 & 10,878 \\
60 & 8 & 8 & 19,071 & 1,830 & 94 & 11 & 5 & 87,878 & 11,418 \\
61 & 9 & 7 & 11,055 & 773 & 95 & 8 & 8 & 86,076 & 8,418 \\
62 & 5 & 11 & 14,345 & 1,116 & 96 & 7 & 9 & 101,613 & 14,413 \\
63 & 9 & 7 & 20,941 & 1,632 & 97 & 1 & 9 & 57,396 & 6,064 \\
64 & 7 & 9 & 21,638 & 1,735 & 98 & 7 & 9 & 94,705 & 13,529 \\
65 & 9 & 7 & 20,678 & 2,335 & 99 & 8 & 8 & 113,528 & 25,340 \\
66 & 8 & 8 & 24,944 & 4,649 & 100 & 8 & 8 & 116,954 & 8,518 \\
\midrule
\multicolumn{10}{l}{all: 606 identity entries and 465 bounds improved, 2,218,809 terms saved} \\
\bottomrule
\end{tabular}

\caption{Improvements by power for $33\le k\le100$ ($m\le16$), by all the methods of
Section~\ref{sec:high}: entries improved that were attained by an identity and that were
derived bounds, terms saved, and the largest saving in one entry.}\label{tab:perkhigh}
\end{table}

\begin{table}[htbp]
\centering\small
% generated by tools/contrib.py -- do not edit
\begin{tabular}{lrrrr}
\toprule
method & improves & only this & cumulative & terms saved \\
\midrule
sieve search & 795 & 688 & 795 & 1,572,286 \\
coset starts & 188 & 98 & 894 & 1,928,740 \\
known-identity starts & 16 & 14 & 908 & 1,935,133 \\
Wr\'oblewski completions & 160 & 132 & 1,040 & 2,089,106 \\
structured completions & 31 & 31 & 1,071 & 2,218,809 \\
\bottomrule
\end{tabular}

\caption{Contribution of each method to the improvements for $33\le k\le100$ ($m\le16$).
``Sieve search'': the searches of Section~\ref{sec:highsearch} with all values of $N$ and all
repetitions. ``Improves'': entries improved by the method; ``only this'': entries improved by it
and by no other method; ``cumulative'' and ``terms saved'': entries improved, and terms saved by
the best identity found, when the methods are added in the order listed.}\label{tab:contrib}
\end{table}

\subsection{A congruence obstruction}\label{sec:congruence}
The powers that resist the search have an arithmetic feature in common: for some small $j$,
here $1\le j\le4$, the number $k-j$ has many divisors $\ell-1$ with $\ell$ prime.

\begin{proposition}\label{prop:moment}
Let $1\le j<k$, and let $M_j(k)$ be the product of the primes $\ell$ with $(\ell-1)\mid(k-j)$. Every
relation $c\in L_{k,N}$ satisfies
\[
\sum_{x=1}^{N} c_x\,x^j\equiv 0 \pmod{M_j(k)}.
\]
If moreover $M_j(k)>\|c\|_1N^j$, then $\sum_x c_x x^j=0$.
\end{proposition}

\begin{proof}
Let $\ell$ be a prime with $(\ell-1)\mid(k-j)$. If $\ell\nmid x$, then $x^{k-j}\equiv1\pmod\ell$
by Fermat's theorem, so $x^k\equiv x^j$; if $\ell\mid x$, both sides are divisible by $\ell$
since $j\ge1$. Hence $\sum c_x x^j\equiv\sum c_x x^k=0\pmod\ell$ for each such $\ell$, and the
first claim follows from the Chinese remainder theorem. The second follows from
$\bigl|\sum c_x x^j\bigr|\le\|c\|_1N^j$.
\end{proof}

For $j=0$ the same argument gives the familiar condition of Section~\ref{sec:history} on the
numbers of terms prime to $\ell$, and for identities of type $(k,2,2)$ the case $j=1$ is proved in
\cite{Asiryan2025}. The values of $\mu_j(c)=\sum_xc_xx^j$ on $L_{k,N}$ are also easy
to describe.

\begin{lemma}\label{lem:image}
For $1\le j<k$, $\mu_j(L_{k,N})=g\Z$ with $g=g_j(k,N)=\gcd_{2\le x\le N}(x^k-x^j)$, and
$M_j(k)\mid g$.
\end{lemma}

\begin{proof}
Since $1^k=1$, a relation $c$ has $c_1=-\sum_{x\ge2}c_xx^k$, so $c=\sum_{x\ge2}c_x(e_x-x^ke_1)$,
and these $N-1$ vectors form a basis of $L_{k,N}$. Their moments are $x^j-x^k$. The divisibility
follows from Proposition~\ref{prop:moment}.
\end{proof}

The modulus $g$ is always a multiple of $M_j(k)$. For $j=1$ it equals $M_1(k)$ for every power
of Table~\ref{tab:congruence}, but for $j\ge2$ higher prime powers also divide every
$x^k-x^j$: for $k=50$ and $j=2$, both $4$ and $9$ divide $x^{50}-x^2$ for every $x$, and
$g=6\,M_2(50)$; similarly $g=36\,M_3(75)$ and $g=120\,M_4(64)$.

The explanation of the powers that resist the search rests on three statements of different
kinds. (i) Proposition~\ref{prop:moment} and Lemma~\ref{lem:image} are exact. (ii) That short
relations lie in the kernel of $\mu_j$ when $M_j(k)$ is large is an empirical observation: the moment of a short relation is far smaller than the
worst-case bound $\|c\|_1N^j$, and when it is smaller than $M_j(k)$ it vanishes. On the rows of a
BKZ-20 basis of $L_{k,100}$, $\mu_j$ vanishes on more than 90\% of the rows for some $j\le12$ for
\ForcedK{} of the powers $33\le k\le100$ (Table~\ref{tab:congruence}) and for
$k\in\{\ForcedMainK\}$ in the main table; we then call the $j$-th moment forced. In each case $k-j$ is one of $16$, $18$, $24$, $30$, $36$, $40$,
$48$, $60$, $66$, $72$, $80$, $84$, $88$, $90$ or $96$, numbers with many divisors $\ell-1$
with $\ell$ prime. On the saved sieve databases of $L_{k,100}$, with $4.2\times10^6$ vectors each, the moment
vanishes on every vector for all of these powers except $k=17$, where $M_1(17)=510$ is small and
it vanishes on 65\% of them. (iii) A descent is confined to a coset of the kernel only under a
further condition. When every move has $\mu_j=0$, as in these databases, a descent never changes
$\mu_j$ and stays in the coset of its start. The database vectors used as starts lie in the
kernel, and so do most rows of the reduced basis; the few rows with $\mu_j\ne0$ lead to other
cosets, and for $k=50$ the plain search reached the coset $\mu_2=-3g$ in this way.

Heuristically, the kernel contains few very unbalanced relations. If $\mu_j(c)=0$ and one side of
$c$ has $m$ terms, the other side has $\sum|c_x|x^j\le mN^j$, hence $n\le mN^j$. For $j=1$ this
bound is restrictive; for $j\ge2$ it is weak, and beyond it we have only the heuristic that a long
side with small $\sum|c_x|x^j$ has mostly small bases, whose $k$-th powers are hard to balance
against those of a short side. A relation with a small side and a long side would then have
$\mu_j$ equal to a nonzero multiple of $g$, roughly $-n$ times the mean of $x^j$ on its long side,
in a coset that a descent from the kernel never reaches. The data agree with this explanation.
The plain search improved \ForcedSieveImp{} of the \ForcedCells{} entries for the \ForcedK{}
powers with a forced moment, and \UnforcedSieveImp{} of the \UnforcedCells{} entries for the
other \UnforcedK. For $k=82$ and $89$ the moment is not forced at $N=60$, where the search reached
$m=1$, and is forced from $N=70$ on, where it found no relation with $m\le40$.

The cure is to start in the right coset. An extended Euclidean algorithm on the moments of a
reduced basis gives a relation with $\mu_j=g$ (Lemma~\ref{lem:image}), which is shortened by
Babai's algorithm \cite{Babai1986} against the rows with $\mu_j=0$. The descent then starts from
relations with $\mu_j=\pm tg$, obtained in the same way, for several $t$ between $1$ and $20$, and
from these starts plus random database vectors; its moves preserve the coset. We call these
coset starts. For the five powers $k\in\{64,76,82,87,89\}$, where nothing else had worked, they
improve every entry with $m\le16$; for example $(89,1)$ falls from 11{,}431 to 3{,}645 and
$(82,1)$ from 9{,}665 to 2{,}664. Coset starts also give the best identities found for all or
nearly all entries with $k\in\{39,41,67,81,91,98\}$, and for $k=25$ in the main table they give
most of the improvements.

Coset starts do not help when the right coset is out of reach with bases at most
$N$: a relation whose long side has $n$ terms has $|\mu_j|\le n\,N^j$, and on average about
$nN^j/(j+1)$, so a coset $\mu_j=-tg$ with $t\ge1$ needs $n\gtrsim g(j+1)/N^j$. For $k=89$ this bound is about $1{,}200$, well below the records, and coset starts
succeed. For $k=75$, where $g=36\,M_3(75)$, and for $k=85$ and $97$ it is about $20{,}000$,
$68{,}000$ and $90{,}000$, far above the records; there only larger bases help, which is what
the completions of the next two subsections provide. The criterion is necessary but not
sufficient: it identifies the cosets that can contain short unbalanced relations, not the starts
from which a descent reaches them.

For $k=50$, where $g=6\,M_2(50)$ and the bound is about $84$, the coset starts reach $m\ge7$
only with $n$ between two and five times the table entries, with bases up to $110$ as well. The
long sides of these relations have root-mean-square base between $30$ and $36$, against $47$ to
$83$ for the coset-start relations of the other powers of Table~\ref{tab:congruence} and $60$ to
$70$ for the records, so that a given coset needs about three to five times as many terms. The
sieve database of $L_{50,100}$ does not differ from those of the other powers in any statistic we
measured, such as the distribution of its mass over the bases or the change in the number of
terms per move. The records for $k=50$ use bases up to $117$ and lie in the cosets $t=10$ and
$t=15$. Two runs of coset starts with bases $1,\dots,117$, one in these two cosets and one in
five others, reach no relation with $m<40$ and $m<27$, respectively; their starts have median
numbers of terms 3,300 and 2,100, two to four times that of the records. For this power the
descent from coset starts does not reach the unbalanced relations even in the cosets that
contain them, and we do not know why this happens for $k=50$ and not for the other powers. Known-identity
starts work instead: from the 22 identities with $m\le40$ in Meyrignac's database, all of which
have bases at most $117$, the descent with the sieve database of $L_{50,117}$ as moves improves
all 16 entries, $(50,1)$ from 1{,}164 to 577. The new relations lie in the cosets $t=5$ and
$t=10$, and their long sides have root-mean-square base between $59$ and $78$.

\begin{table}[htbp]
\centering\scriptsize
\setlength{\tabcolsep}{4pt}
% generated by tools/congruence.py -- do not edit
\begin{tabular}{rrlrrr}
\toprule
$k$ & $j$ & primes $\ell$ with $(\ell-1)\mid(k-j)$ & vanishes & plain & all \\
\midrule
37 & 1 & 2, 3, 5, 7, 13, 19, 37 & 99\% & 0/16 & 16/16 \\
38 & 2 & 2, 3, 5, 7, 13, 19, 37 & 99\% & 0/16 & 16/16 \\
39 & 3 & 2, 3, 5, 7, 13, 19, 37 & 99\% & 14/16 & 16/16 \\
40 & 4 & 2, 3, 5, 7, 13, 19, 37 & 97\% & 0/16 & 16/16 \\
41 & 1 & 2, 3, 5, 11, 41 & 99\% & 12/16 & 16/16 \\
42 & 2 & 2, 3, 5, 11, 41 & 97\% & 0/16 & 16/16 \\
49 & 1 & 2, 3, 5, 7, 13, 17 & 99\% & 11/16 & 16/16 \\
50 & 2 & 2, 3, 5, 7, 13, 17 & 98\% & 2/16 & 16/16 \\
61 & 1 & 2, 3, 5, 7, 11, 13, 31, 61 & 99\% & 0/16 & 16/16 \\
62 & 2 & 2, 3, 5, 7, 11, 13, 31, 61 & 99\% & 0/16 & 16/16 \\
63 & 3 & 2, 3, 5, 7, 11, 13, 31, 61 & 99\% & 0/16 & 16/16 \\
64 & 4 & 2, 3, 5, 7, 11, 13, 31, 61 & 98\% & 0/16 & 16/16 \\
67 & 1 & 2, 3, 7, 23, 67 & 99\% & 12/16 & 16/16 \\
73 & 1 & 2, 3, 5, 7, 13, 19, 37, 73 & 99\% & 0/16 & 16/16 \\
74 & 2 & 2, 3, 5, 7, 13, 19, 37, 73 & 99\% & 0/16 & 16/16 \\
75 & 3 & 2, 3, 5, 7, 13, 19, 37, 73 & 99\% & 0/16 & 12/16 \\
76 & 4 & 2, 3, 5, 7, 13, 19, 37, 73 & 99\% & 0/16 & 16/16 \\
81 & 1 & 2, 3, 5, 11, 17, 41 & 99\% & 14/16 & 16/16 \\
82 & 2 & 2, 3, 5, 11, 17, 41 & 99\% & 0/16 & 16/16 \\
85 & 1 & 2, 3, 5, 7, 13, 29, 43 & 99\% & 0/16 & 9/16 \\
86 & 2 & 2, 3, 5, 7, 13, 29, 43 & 99\% & 0/16 & 16/16 \\
87 & 3 & 2, 3, 5, 7, 13, 29, 43 & 99\% & 0/16 & 16/16 \\
89 & 1 & 2, 3, 5, 23, 89 & 99\% & 0/16 & 16/16 \\
91 & 1 & 2, 3, 7, 11, 19, 31 & 99\% & 12/16 & 16/16 \\
97 & 1 & 2, 3, 5, 7, 13, 17, 97 & 99\% & 0/16 & 10/16 \\
98 & 2 & 2, 3, 5, 7, 13, 17, 97 & 99\% & 14/16 & 16/16 \\
\bottomrule
\end{tabular}

\caption{Powers $33\le k\le100$ whose short relations have a forced moment: the degree $j$, the
primes in $M_j(k)$, the fraction of the rows of a BKZ-20 basis of $L_{k,100}$ on which $\mu_j$
vanishes, and the entries ($m\le16$) improved by the plain search with bases $1,\dots,100$
(``plain'') and by all methods (``all''), including the coset and known-identity starts of
Section~\ref{sec:congruence} and the completions of Sections~\ref{sec:completion}
and~\ref{sec:structured}.}\label{tab:congruence}
\end{table}

The proposition also explains why the method that produced the existing records for these
powers works. In Wr\'oblewski's congruence searches and the completions of
Section~\ref{sec:history} that the author ran on them in 2002, the congruence part supplies a
large nonzero moment, which places the completed relation in a coset where unbalanced relations
exist. Many records in the table are of this kind, for example those for $k=37$, $38$, $61$,
$73$ and $97$.

\subsection{Completing congruence solutions}\label{sec:completion}
The completion of 2002 (Section~\ref{sec:history}) starts from bases $y_1,\dots,y_r$, found by
Wr\'oblewski, with $\sum y_i^k$ divisible by $M^k$ for a large $M$, and writes
$\Delta=\sum y_i^k/M^k$ as $\sum_{x\le N}c_x x^k$. This gives the identity
\[
\sum_{c_x>0}c_x\,(Mx)^k=\sum_{i=1}^{r}y_i^k+\sum_{c_x<0}|c_x|\,(Mx)^k .
\]
The congruence solutions supply the part of the relation that the moment conditions of
Proposition~\ref{prop:moment} make hard to reach, and the lattice of small bases completes it.
The 2003 archive contains Wr\'oblewski's congruence solutions for
$k=31$, $37$, $38$, $49$, $61$, $62$, $63$, $68$, $73$, $74$, $85$ and $86$, as well as for
$k=19$ and $26$; we call them the historical inputs. Most of them are one-sided, with all bases
$y_i$ on one side as above;
the signed ones have bases on both sides, as $P$ and $Q$ in Section~\ref{sec:history}. We
repeated the completion with the methods of Section~\ref{sec:method}. A particular solution
$c^{(0)}$ comes from a reduced row with last coordinate $0$ and middle coordinate $\pm1$ of the
embedding basis
$\bigl(\begin{smallmatrix}I&0&\lambda v\\0&1&-\lambda\Delta\end{smallmatrix}\bigr)$ after LLL and
BKZ-20; the extra coordinate, which the variant of 2002 did not have, identifies the multiple of
$\Delta$. The descent of Algorithm~\ref{alg:descent} then runs on $c$ with the sieve database of
$L_{k,100}$ (later also $L_{k,110}$) as moves, with the fixed terms $y_i$ added to the count of
their side, and starts at $c^{(0)}$ and at $c^{(0)}\pm u$ for 100 random database vectors $u$. A
fixed base $y_i=Mx$ with $x\le N$ is first absorbed into $c_x$, so that the remaining fixed bases
differ from every base $Mx$ of the completion and the descent counts the terms exactly. Every
identity is assembled and verified in exact arithmetic.

With bases at most $100$ the completions improve every entry with $m\le16$ for
$k\in\{37,38,49,61,62,63,73,74\}$ and 15 of the 16 for $k=86$. For eight of these
nine powers the plain search of Section~\ref{sec:highsearch} had improved no entry, and for
$k=49$ it had improved 11 (Table~\ref{tab:congruence}, column ``plain''). For $k=68$ the completions
improve every entry but do not beat the sieve search, and for $k=85$ they improve none. For
$k=31$ of the main table, where the records use bases from 127 to 265 and no sieve search
improved an entry, the completions improve all \KThirtyOneImproved{} entries, by
\KThirtyOneSaved{} terms in all. For this power we also completed the signed inputs, and the
one-sided ones with bases up to $140$, using the database of $L_{31,100}$ padded with zeros
together with the rows of 100 BKZ-30 reductions of $L_{31,140}$. The 23{,}746 one-sided and 80{,}345 signed inputs for $k=19$ and $k=26$ gave
identities below the table for five entries with $k=26$, but none shorter than those of the
sieve searches; the one improvement on our own results, of type
$(26,14,60)$, is still above the table entry $52$.

\subsection{Structured congruence solutions}\label{sec:structured}
For the powers without historical inputs we constructed congruence solutions ourselves. Two
observations make this work. The first is structural. Wr\'oblewski's inputs, sorted by the
$p$-adic valuation of their bases, fall into levels: his input for $k=37$ has 4, 5, 6, 5, 7 and 11
bases of valuations $0,\dots,5$ for $M=2^6$. With bases $y=p^is$, $p\nmid s$, $0\le i<a$ and
$M=p^a$, write $A_i=\sum_s u_s s^k$ for level $i$, where $u_s$ is the signed multiplicity of the
base $p^is$. Then $\sum_y u_yy^k=\sum_ip^{ki}A_i$, and divisibility by $M^k$ is the recursion
$h_0=0$, $h_{i+1}=(h_i+A_i)/p^k$ for the carries $h_i$, with $\Delta=h_a$: level $i$ needs only
$A_i\equiv-h_i\pmod{p^k}$, a congruence in a lattice of determinant $p^k$ instead of $p^{ka}$.
Each level is solved by a short affine start and a descent inside its lattice that favors one
sign. We used $p=2$, and for even $k$ also $p=3$: for $p=2$ and even $k$, every odd $s$ satisfies
$s^k\equiv1\pmod{2^{v_2(k)+2}}$, which fixes the number of terms of a one-sided level modulo
$2^{v_2(k)+2}$.

The second observation is a condition on the coset of the completion. Since $g=g_j(k,N)$
divides $x^k-x^j$ for every $x\le N$ (Lemma~\ref{lem:image}), every completion $c$ with
$\sum_xc_xx^k=\Delta$ satisfies $\sum_xc_xx^j\equiv\Delta\pmod g$, and the $j$-th moments of the
completions are exactly the integers $\equiv\Delta\pmod g$. A completion leans to one side, as
the records with small $m$ require, only if its moment is a moderate negative number, so the
residue of $\Delta$ modulo $g$, taken in $(-g/2,g/2]$, should lie in a window $W$ of such values.
This is the remark in a message of Wr\'oblewski of 14 October 2002 (personal communication) that
a number to be decomposed into 31st powers ``has remainder 7017 mod 14322'', where
$14322=M_1(31)=g_1(31,N)$: the first powers of its terms must sum to an integer
$\equiv7017\pmod{14322}$. We took $W=[-\tfrac N2n_1,-\tfrac N{16}n_1]$, with $n_1$ the table entry
for $m=1$, the range of the first moment of a completion with $n_1/8$ to $n_1$ terms of mean
base $N/2$. Our implementation tested the residue of $\Delta$ modulo $M'$, the odd part of
$M_j(k)$, instead of modulo $g$; here $g=2M'$ for $k=85$ and $97$, and $g=72M'$ for $k=75$. This
test is weaker: for $k=85$ and $97$ it also keeps solutions whose moments are $r\pm M'$,
$r\pm3M'$, \dots{} with $r\in W$, none of them in $W$. Of the 30 solutions kept for each of these
two powers, 15 have $\Delta\bmod g$ in $W$, and all the improvements below come from these. For
$k=75$, where $j=3$, the window is not matched to the third moment, and none of the 11 solutions
kept has $\Delta\bmod g$ in $W$; the completions of one of them still produced the improvements
reported below. Algorithm~\ref{alg:struct} summarizes the construction as it was run, and
Table~\ref{tab:struct} gives its parameters.

\begin{algorithm}[ht]
\caption{Structured congruence solutions}\label{alg:struct}
\begin{algorithmic}[1]
\Require $k$, prime $p$, $M=p^a$, candidates per level $C$, bound $B$, modulus $M'$, window $W$
\Repeat
  \State $h\gets0$; $P,Q\gets\emptyset$
  \For{$i=0,\dots,a-1$}
    \State draw $C$ integers $s\le\max(B/p^i,2C+2)$ with $p\nmid s$, at random with weights
    increasing linearly in $s$
    \State reduce and sieve $\Lambda=\{u\in\Z^C:\sum_su_ss^k\equiv0\pmod{p^k}\}$
    \State $u_0\gets$ a one-term solution of $\sum_su_ss^k\equiv-h\pmod{p^k}$, reduced by Babai
    against $\Lambda$
    \State descend from $u_0$ and $u_0\pm$ random vectors of $\Lambda$, moves $\Lambda$; keep the
    $u$ minimizing $\|u\|_1+4\sum_{u_s<0}|u_s|$
    \State $h\gets(h+\sum_su_ss^k)/p^k$; add $|u_s|$ copies of $p^is$ to $P$ if $u_s>0$, to $Q$ if
    $u_s<0$
  \EndFor
  \If{$h\bmod M'$, taken in $(-M'/2,M'/2]$, lies in $W$} \textbf{output} $(P,Q)$, with
  $\sum_{y\in P}y^k-\sum_{y\in Q}y^k=h\,M^k$ \EndIf
\Until{enough solutions or tries}
\end{algorithmic}
\end{algorithm}

\begin{table}[htbp]
\centering\small
\begin{tabular}{rrrrrrrr}
\toprule
$k$ & $M$ & $C$ & $B$ & $M'$ & window $W$ & tries & kept \\
\midrule
75 & $2^6$ & 48 & 2{,}048 & 70{,}050{,}435 & $[-408{,}000,-51{,}000]$ & 2{,}000 & 11 \\
85 & $2^7$ & 48 & 4{,}096 & 1{,}702{,}155 & $[-294{,}650,-36{,}832]$ & 145 & 30 \\
97 & $2^6$ & 48 & 2{,}048 & 2{,}250{,}885 & $[-568{,}100,-71{,}013]$ & 119 & 30 \\
\bottomrule
\end{tabular}
\caption{Parameters of the structured congruence solutions whose completions gave the records
for $k=75$, $85$ and $97$ ($p=2$, $B=32M$): the modulus $M=p^a$, candidates per level $C$,
bound $B$ on the bases of level $0$, the modulus $M'$ and window $W$ for $\Delta$, and the numbers of
tries and of solutions kept.}\label{tab:struct}
\end{table}

These congruence solutions have between about 200 and 360 bases, nearly all on one side.
Larger moduli, up to $2^9$, gave no further improvement. Completed as above, with bases at most
$100$ and, in the repetitions of Section~\ref{sec:variation}, $110$, they improve entries for
three powers that no other method had touched: those with $m\ge5$ for $k=75$ (for example $(75,6)$ from 6{,}266 to
1{,}991), those with $m\ge8$ for $k=85$, and those with $m\ge7$ for $k=97$. The entries with
smaller $m$ for these three powers are the only ones in the second table that we did not
improve. The structured solutions gave nothing for $k\in\{64,76,82,87,89\}$, which the coset starts of
Section~\ref{sec:congruence} later resolved completely.

%======================================================================
\section{Discussion}\label{sec:discussion}
%======================================================================

Two kinds of conclusion should be kept apart. Section~\ref{sec:experiments} evaluates methods
under controlled budgets on a few lattices, with three seeds each; Sections~\ref{sec:results}
and~\ref{sec:high} report the best identities found by many searches that differ in $N$, $\rho$,
method and number of repetitions, and so measure record discovery, not the expected yield of a
method. Within these limits the experiments suggest the following. The records are not short
vectors: the shortest relations of $L_{k,N}$ are balanced, and for $20\le k\le30$ a database of
half a million to a million of them at $N=60$ or $70$ contains almost no relation with
$\min(p,q)\le16$. A descent turns balanced relations into unbalanced ones in a few steps (the
median is \StepsMedian), provided the move set is large; the 2003 procedure, whose moves are the
rows of one reduced basis, is far behind except for $k=20$. The size of the move set matters
most: at equal size neither BKZ-40 rows nor sieve vectors are uniformly better, but the sieve
supplies many more moves in the same time. The records with small $m$ are found at small
$N$, far below the dimension that the Gaussian heuristic suggests. Combining known identities
with one another adds nothing here, but known identities are useful as starts, especially at
small $N$. Where the search
fails for the higher powers, the cause is usually arithmetic: a low moment vanishes on almost all
short relations, a descent from such a relation cannot leave its coset, and for most of these
powers starting in another coset removes the obstruction. The searches are randomized and the reported values are not
converged: each round of repetitions lowered hundreds of entries.

The results have clear limits. We make no claim of minimality; the descent is greedy and
offers no completeness guarantee. The searches use bases at most $110$ for most powers and
at most $120$ in all, and only the completions of Section~\ref{sec:high} use larger bases,
through the congruence part and the factor $M$. The guidance of Table~\ref{tab:gh} is heuristic for this structured lattice. The
reach of the method is limited by the size of the sieve database, which determines both the
memory needed and the cost of the descent. We did not attempt the table of prime bases.

Several questions remain open.
\begin{enumerate}
\item Sieving minimizes $\|c\|_2$, while a record is judged by $\min(p,q)$ and then by
$\|c\|_1$, and the descent connects the two greedily. A sieve or enumeration that targets these
quantities directly, for example with one-sided pruning, might find the unbalanced relations
without the descent.
\item Table~\ref{tab:gh} suggests that $N$ between about 210 and 300 would allow shorter
balanced relations for $24\le k\le32$. A full sieve at that size is beyond current memory, so a
method with smaller memory requirements, such as enumeration or tuple sieving
\cite{BaiLaarhovenStehle2016}, or a form of dimensions for free suited to this objective, would
be needed.
\item For $k=75$, $85$ and $97$ the entries with small $m$ need a coset that bases at most $110$
cannot reach, and the structured congruence solutions of Section~\ref{sec:structured} have so
far reached only the entries with $m\ge5$, $m\ge8$ and $m\ge7$, respectively. Congruence solutions with fewer bases, or completions over a larger
range of bases, might reach the remaining entries.
\item The unimproved entries with small bases for $17\le k\le21$, most of them for $k=18$, and
the failure of the coset starts for $k=50$ are not explained by Proposition~\ref{prop:moment}
alone.
\end{enumerate}

\section*{Acknowledgments}
We thank Jean-Charles Meyrignac for maintaining the record tables for more than two decades;
the volunteers of EulerNet and yoyo@home; Uwe Beckert and Robert Gerbicz; Jaros\l aw
Wr\'oblewski, whose congruence solutions of 2002 are used in Section~\ref{sec:completion}; and
the contributors
named in Section~\ref{sec:history} and Appendix~\ref{app:chronology}. This work used the Turing
cluster of the Titan Supercomputing Center at California State University, Fullerton (CSUF). This work
used resources available through the National Research Platform (NRP) at the University of
California, San Diego. NRP has been developed, and is supported in part, by funding from
National Science Foundation, from awards 1730158, 1540112, 1541349, 1826967, 2112167, 2100237,
and 2120019, as well as additional funding from community partners. The Titan Supercomputing
Center at CSUF is one of the partners contributing to NRP resources.

\medskip\noindent\textit{Use of generative AI.}
The author used Claude (Anthropic), a large language model, to assist with porting the original
software, developing and running the search programs, searching the literature, and
preparing the manuscript. The author designed and directed the work and takes full
responsibility for its content. All identities reported were verified independently in exact
integer arithmetic, and every reference was checked against Crossref, zbMATH Open, arXiv, or the
current or archived version of the cited web page.

\medskip\noindent\textit{Data availability.}
The code and data are in the repository \url{https://github.com/gchilders/sumsbkz-py}, archived
at Zenodo as \href{https://doi.org/10.5281/zenodo.23849788}{\nolinkurl{doi:10.5281/zenodo.23849788}}. Its
data archive (directory \texttt{archive}) contains the identities, the reference values of the
table, the standalone verification program, and SHA-256 checksums of the snapshot of
Meyrignac's pages and database from which the reference values were extracted. The repository
also contains the search code, the rebuilt \texttt{SumsBKZ}, the record databases of 2001--2003,
and the scripts that produce every table and number in this paper. The programs \texttt{NewLC}, \texttt{ResultLC} and \texttt{ysum}, written by others,
are described from the author's archive and are not distributed, and neither are
Wr\'oblewski's congruence solutions; the identities obtained from them are in the archive.

\appendix

\section{Chronology of the distributed searches}\label{app:chronology}
Meyrignac's distributed search began on 1 January 1999, with ranges reserved by e-mail; the
clients \texttt{Sum95} and \texttt{Sum99} followed during 1999. The client \texttt{Euler2000} and the automated EulerNet
client--server system followed in July 2000, with the server hosted by the author from the
start. Hosting moved to Nuutti Kuosa in December 2001 and returned to the author in 2002, and
the server ran until 2010. The search moved from $(6,1,7)$ to $(6,1,6)$ in August 2000 and to
$(6,2,5)$ in March 2002. Its last statistics, from October 2010, record about 309 CPU-years
\cite{MeyrignacEuler}, and by then all solutions of $(6,2,5)$ with largest term below $88{,}000$ were known
\cite{GerbiczMeyrignacBeckert2011}. Resta and Meyrignac determined the smallest solutions of
$x^6+y^6=a^6+b^6+c^6+d^6+e^6$ \cite{RestaMeyrignac2003}.

From April 2010 to July 2011 the $(6,2,5)$ enumeration continued as the Euler~(6,2,5)
subproject of yoyo@home, a BOINC project \cite{Anderson2004,Anderson2020} provided and
operated by Uwe Beckert, with an algorithm of Robert Gerbicz
\cite{YoyoEuler625,MeyrignacEuler625Forum2010,GerbiczMeyrignacBeckert2011}. It ran in two phases, to $7^6=117{,}649$ (about
70 GHz-years) and then to $250{,}000$ (about 740 GHz-years in all), and established that there are
exactly 377 primitive solutions with all terms below $250{,}000$; 4,470 volunteers computed at
least one work unit \cite{GerbiczMeyrignacBeckert2011}.

\clearpage
\begin{landscape}
\section{All entries for \texorpdfstring{$17\le k\le32$}{17 <= k <= 32}}\label{app:grid}
\begin{center}
\scriptsize
% generated by tools/make_tables.py -- do not edit
\setlength{\tabcolsep}{1.6pt}
\begin{tabular}{rcccccccccccccccc}
\toprule
$k\backslash m$ & 1 & 2 & 3 & 4 & 5 & 6 & 7 & 8 & 9 & 10 & 11 & 12 & 13 & 14 & 15 & 16 \\
\midrule
17 & \textcolor{gray}{39} & \textbf{37}/40 & \textbf{34}/35 & \textbf{27}/33 & \textbf{26}/27 & \textbf{26}/28 & \textcolor{gray}{23} & \textbf{22}/24 & \textbf{21}/23 & \textcolor{gray}{20} & \textbf{18}/20 & \textbf{16}/18 & \textcolor{gray}{16} & \textcolor{gray}{14} &  &  \\
18 & \textcolor{gray}{57} & \textcolor{gray}{57} & \textbf{42}/57 & \textbf{41}/44 & \textbf{41}/43 & \textcolor{gray}{34} & \textcolor{gray}{35} & \textbf{34}/36 & \textcolor{gray}{31} & \textcolor{gray}{28} & \textcolor{gray}{29} & \textcolor{gray}{25} & \textcolor{gray}{15} & \textcolor{gray}{16} & \textcolor{gray}{15} &  \\
19 & \textcolor{gray}{51} & \textcolor{gray}{52} & \textbf{49}/53 & \textbf{46}/51 & \textbf{42}/43 & \textbf{39}/40 & \textbf{36}/41 & \textbf{30}/36 & \textbf{28}/33 & \textbf{26}/29 & \textbf{26}/30 & \textbf{25}/31 & \textbf{21}/23 & \textbf{18}/24 & \textcolor{gray}{17} & \textbf{16}/17 \\
20 & \textcolor{gray}{61} & \textbf{56}/62 & \textbf{56}/61 & \textbf{54}/60 & \textbf{51}/59 & \textbf{47}/53 & \textbf{37}/39 & \textbf{38}/40 & \textcolor{gray}{35} & \textcolor{gray}{35} & \textbf{33}/35 & \textbf{34}/35 & \textbf{33}/35 & \textbf{25}/35 & \textbf{27}/35 & \textbf{18}/26 \\
21 & \textbf{67}/75 & \textbf{63}/72 & \textbf{55}/67 & \textbf{51}/60 & \textbf{42}/57 & \textbf{41}/50 & \textbf{35}/51 & \textbf{33}/46 & \textbf{34}/41 & \textbf{32}/39 & \textbf{27}/37 & \textbf{28}/34 & \textbf{23}/26 & \textcolor{gray}{25} & \textbf{23}/25 & \textbf{20}/24 \\
22 & \textbf{71}/95 & \textbf{73}/75 & \textbf{69}/76 & \textbf{70}/72 & \textbf{51}/73 & \textbf{51}/73 & \textbf{51}/59 & \textbf{52}/57 & \textbf{52}/54 & \textbf{52}/54 & \textbf{37}/52 & \textbf{35}/36 & \textbf{36}/37 & \textbf{35}/37 & \textbf{35}/36 & \textbf{35}/37 \\
23 & \textbf{82}/105 & \textbf{76}/91 & \textbf{64}/87 & \textbf{66}/86 & \textbf{63}/81 & \textbf{59}/72 & \textbf{47}/62 & \textbf{42}/63 & \textbf{40}/62 & \textbf{38}/53 & \textbf{32}/49 & \textbf{32}/47 & \textbf{33}/41 & \textbf{32}/36 & \textbf{30}/37 & \textbf{29}/33 \\
24 & \textbf{104}/124 & \textbf{97}/116 & \textbf{85}/109 & \textbf{73}/97 & \textbf{74}/85 & \textbf{73}/85 & \textbf{62}/71 & \textbf{62}/72 & \textbf{61}/65 & \textbf{51}/66 & \textbf{56}/57 & \textbf{52}/58 & \textbf{51}/58 & \textbf{45}/55 & \textbf{39}/52 & \textbf{40}/53 \\
25 & \textbf{111}/137 & \textbf{93}/118 & \textbf{89}/104 & \textbf{85}/94 & \textbf{82}/89 & \textbf{80}/90 & \textbf{79}/81 & \textbf{75}/82 & \textbf{68}/80 & \textbf{65}/80 & \textbf{65}/74 & \textbf{63}/70 & \textbf{59}/68 & \textbf{59}/69 & \textbf{57}/69 & \textbf{57}/67 \\
26 & \textbf{135}/155 & \textbf{110}/133 & \textbf{101}/119 & \textbf{102}/116 & \textbf{90}/111 & \textbf{86}/106 & \textbf{73}/99 & \textbf{69}/100 & \textbf{73}/81 & \textbf{68}/82 & \textbf{67}/75 & \textbf{62}/72 & \textbf{60}/70 & \textcolor{gray}{52} & \textbf{47}/53 & \textbf{48}/54 \\
27 & \textbf{117}/162 & \textbf{108}/146 & \textbf{91}/132 & \textbf{86}/118 & \textbf{84}/119 & \textbf{79}/104 & \textbf{81}/99 & \textbf{71}/100 & \textbf{65}/87 & \textbf{64}/88 & \textbf{59}/89 & \textbf{52}/63 & \textbf{50}/64 & \textbf{46}/65 & \textbf{46}/66 & \textbf{40}/56 \\
28 & \textbf{126}/183 & \textbf{119}/147 & \textbf{118}/148 & \textbf{95}/128 & \textbf{93}/129 & \textbf{93}/124 & \textbf{92}/124 & \textbf{92}/124 & \textbf{93}/96 & \textbf{70}/96 & \textbf{69}/96 & \textbf{68}/97 & \textbf{70}/98 & \textbf{70}/74 & \textbf{48}/73 & \textbf{46}/74 \\
29 & \textbf{132}/173 & \textbf{130}/168 & \textbf{118}/145 & \textbf{108}/146 & \textbf{104}/147 & \textbf{90}/136 & \textbf{87}/137 & \textbf{93}/119 & \textbf{86}/120 & \textbf{82}/108 & \textbf{74}/109 & \textbf{75}/102 & \textbf{64}/96 & \textbf{60}/89 & \textbf{57}/90 & \textbf{50}/84 \\
30 & \textbf{157}/191 & \textbf{138}/188 & \textbf{131}/189 & \textbf{128}/164 & \textbf{129}/139 & \textbf{101}/140 & \textbf{101}/141 & \textbf{100}/142 & \textbf{102}/133 & \textbf{102}/134 & \textbf{79}/135 & \textbf{82}/106 & \textbf{78}/107 & \textbf{76}/108 & \textbf{77}/109 & \textbf{50}/80 \\
31 & \textbf{163}/211 & \textbf{149}/200 & \textbf{136}/199 & \textbf{128}/191 & \textbf{125}/184 & \textbf{124}/182 & \textbf{121}/160 & \textbf{117}/161 & \textbf{113}/162 & \textbf{104}/162 & \textbf{99}/153 & \textbf{87}/151 & \textbf{94}/149 & \textbf{91}/150 & \textbf{85}/143 & \textbf{85}/127 \\
32 & \textbf{194}/230 & \textbf{203}/214 & \textbf{183}/210 & \textbf{179}/210 & \textbf{181}/194 & \textbf{178}/195 & \textbf{170}/191 & \textbf{167}/179 & \textbf{167}/180 & \textbf{165}/181 & \textbf{163}/179 & \textbf{158}/175 & \textbf{155}/168 & \textbf{157}/167 & \textbf{156}/157 & \textbf{157}/158 \\
\bottomrule
\end{tabular}

\captionof{table}{Entries with $m\le16$ of Meyrignac's table for $17\le k\le32$, accessed 27
October 2026. Improved entries show the new $n$ in bold and the reference value after the
slash; other entries show the reference value in gray. Blank entries are empty in the table;
the entry $(18,15)$, also empty in the table, shows the value of an identity in its database. The
counts of Section~\ref{sec:results} refer to $17\le k\le30$; the entries for $k=31$ and $32$
are discussed in Sections~\ref{sec:notimproved} and~\ref{sec:completion}.}\label{tab:grid}
\end{center}
\end{landscape}

\bibliographystyle{amsplain}
\bibliography{literature/references}

\end{document}
